\documentclass[10pt,letterpaper, openany]{book}

% ---------- Idioma y codificación ----------
\usepackage[utf8]{inputenc}
\usepackage[T1]{fontenc}
\usepackage[spanish]{babel}

% ---------- Matemáticas ----------
\usepackage{amsmath}
\usepackage{amssymb}
\usepackage{amsthm}
\usepackage{mathtools}
\usepackage{cancel}

% ---------- Geometría de página ----------
\usepackage[margin=2cm]{geometry}

% ---------- Cajas de color para definiciones/teoremas ----------
\usepackage[most]{tcolorbox}
\tcbuselibrary{skins,breakable}

% ---------- Gráficos ----------
\usepackage{tikz}
\usetikzlibrary{arrows.meta,calc,decorations.markings,decorations.pathreplacing,patterns}

% ---------- Otros ----------
\usepackage{enumitem}
\usepackage{hyperref}
\hypersetup{
    colorlinks=true,
    linkcolor=blue!50!black,
    citecolor=blue!50!black,
    urlcolor=blue!50!black
}
\usepackage{xcolor}

% ================= Entornos tipo teorema =================
\theoremstyle{plain}
\newtheorem{teorema}{Teorema}[section]
\newtheorem{proposicion}[teorema]{Proposición}
\newtheorem{axioma}[teorema]{Axioma}
\newtheorem{corolario}[teorema]{Corolario}

\theoremstyle{definition}
\newtheorem{definicionthm}[teorema]{Definición}
\newtheorem{ejemplothm}[teorema]{Ejemplo}
\newtheorem{observacionthm}[teorema]{Observación}

% Colores
\definecolor{colordef}{RGB}{31,81,144}
\definecolor{colorej}{RGB}{46,125,50}
\definecolor{colorprop}{RGB}{136,14,79}
\definecolor{colorobs}{RGB}{110,110,110}
\definecolor{colorax}{RGB}{191,54,12}

% Cajas estilo tcolorbox que envuelven los teoremas de amsthm
\newtcolorbox{defbox}{
  breakable, enhanced, colback=colordef!5, colframe=colordef,
  fonttitle=\bfseries, boxrule=0.6pt, arc=1.5mm,
  left=6pt, right=6pt, top=4pt, bottom=4pt
}
\newtcolorbox{ejbox}{
  breakable, enhanced, colback=colorej!5, colframe=colorej,
  fonttitle=\bfseries, boxrule=0.6pt, arc=1.5mm,
  left=6pt, right=6pt, top=4pt, bottom=4pt
}
\newtcolorbox{propbox}{
  breakable, enhanced, colback=colorprop!5, colframe=colorprop,
  fonttitle=\bfseries, boxrule=0.6pt, arc=1.5mm,
  left=6pt, right=6pt, top=4pt, bottom=4pt
}
\newtcolorbox{obsbox}{
  breakable, enhanced, colback=colorobs!5, colframe=colorobs,
  fonttitle=\bfseries, boxrule=0.6pt, arc=1.5mm,
  left=6pt, right=6pt, top=4pt, bottom=4pt
}
\newtcolorbox{axbox}{
  breakable, enhanced, colback=colorax!5, colframe=colorax,
  fonttitle=\bfseries, boxrule=0.6pt, arc=1.5mm,
  left=6pt, right=6pt, top=4pt, bottom=4pt
}

% Entornos "amigables" que combinan amsthm + tcolorbox
\newenvironment{definicion}[1][]{%
  \begin{defbox}\begin{definicionthm}[#1]}{\end{definicionthm}\end{defbox}}
\newenvironment{ejemplo}[1][]{%
  \begin{ejbox}\begin{ejemplothm}[#1]}{\end{ejemplothm}\end{ejbox}}
\newenvironment{proposicionc}[1][]{%
  \begin{propbox}\begin{proposicion}[#1]}{\end{proposicion}\end{propbox}}
\newenvironment{teoremac}[1][]{%
  \begin{propbox}\begin{teorema}[#1]}{\end{teorema}\end{propbox}}
\newenvironment{observacion}[1][]{%
  \begin{obsbox}\begin{observacionthm}[#1]}{\end{observacionthm}\end{obsbox}}
\newenvironment{axiomac}[1][]{%
  \begin{axbox}\begin{axioma}[#1]}{\end{axioma}\end{axbox}}
  % en el preámbulo, junto a los demás entornos
\renewcommand{\vec}[1]{\overrightarrow{#1}}
\newcommand{\ihat}{\hat{\imath}}
\newcommand{\jhat}{\hat{\jmath}}
\newcommand{\khat}{\hat{k}}

% ---------- Datos del documento ----------
\pagestyle{plain}
\setlength{\parindent}{0pt}

\title{\textbf{\Huge Mecánica Clásica II}}
\author{}
\date{}

\begin{document}

\frontmatter
\maketitle
\tableofcontents

\mainmatter

\chapter{Mapeos}

\section{Definición y ejemplos básicos}

\begin{definicion}[Mapeo]
Sean $X$ e $Y$ dos conjuntos. Un \textbf{mapeo} $f$ es una regla por la que se
asigna un $y \in Y$ a cada valor $x \in X$. El mapeo se escribe
\[
f : X \to Y .
\]
\end{definicion}

\begin{center}
\begin{tikzpicture}[scale=1.1]
  \draw (0,0) ellipse (1.1 and 1.6);
  \draw (3.5,0) ellipse (1.1 and 1.6);
  \node at (0,1.9) {$X$};
  \node at (3.5,1.9) {$Y$};
  \fill (0,0.3) circle (1.5pt) node[left=2pt] {$x$};
  \fill (3.5,0.3) circle (1.5pt) node[right=2pt] {$y$};
  \draw[-Stealth, thick] (0.3,0.3) .. controls (1.7,1.1) and (1.9,1.1) .. (3.2,0.35);
  \node at (1.75,0.9) {$f$};
\end{tikzpicture}
\end{center}

Cuando $f$ está definido por alguna fórmula explícita, podemos escribir la
correspondencia elemento a elemento como
\[
x \mapsto f(x),
\]
de modo que el mapeo completo se escribe
\[
\begin{aligned}
f : X &\longrightarrow Y \\
x &\longmapsto f(x).
\end{aligned}
\]

\begin{ejemplo}
Sean los conjuntos $X = \{a,b,c\}$ e $Y = \{2,3\}$. Definimos el mapeo
\[
f(a) = 2, \qquad f(b) = 2, \qquad f(c) = 3 .
\]
En este mapeo no hay una fórmula explícita de la aplicación; además, hay dos
elementos de $X$ que corresponden al mismo elemento de $Y$.
\end{ejemplo}

\begin{ejemplo}
Sea el mapeo
\[
\begin{aligned}
f : \mathbb{R} &\to \mathbb{R}\\
x &\mapsto x^2 .
\end{aligned}
\]
Es decir, la función le asocia a cada número real su cuadrado,
$f(x) = x^2$.
\end{ejemplo}

\begin{ejemplo}
Sea el mapeo
\[
\begin{aligned}
f : \mathbb{R}^2 &\to \mathbb{R} \\
(x,y) &\mapsto x + 2y .
\end{aligned}
\]
Este mapeo lleva un punto de $\mathbb{R}^2$ a un punto de $\mathbb{R}$, es
decir, se mapea de dos dimensiones a una dimensión. Un ejemplo físico de este
tipo de mapeo es la temperatura (un número en $\mathbb{R}$) de una lámina
delgada (un objeto bidimensional).
\end{ejemplo}

\section{Imagen inversa, dominio y rango}

\begin{definicion}[Imagen inversa]
Un subconjunto de $X$ cuyos elementos se asignan a $y \in Y$ bajo $f$ se
denomina \textbf{imagen inversa} de $y$, y se denota por
\[
f^{-1}(y) = \{x \in X \mid f(x) = y\}.
\]
\end{definicion}

\begin{center}
\begin{tikzpicture}[scale=1.1]
  \draw (0,0) ellipse (1.2 and 1.7);
  \draw (3.6,0) ellipse (1.1 and 1.6);
  \node at (0,2) {$X$};
  \node at (3.6,1.9) {$Y$};
  \draw[fill=colordef!15] (-0.3,-0.3) circle (0.5);
  \node at (-0.3,-0.3) {\footnotesize$f^{-1}(y)$};
  \fill (3.6,0.3) circle (1.5pt) node[right=2pt] {$y$};
  \draw[-Stealth, thick] (0.2,-0.1) .. controls (1.7,0.9) and (2,0.9) .. (3.3,0.35);
  \node at (1.7,1) {$f$};
\end{tikzpicture}
\end{center}

\begin{ejemplo}
Sea el mapeo
\[
\begin{aligned}
f : \mathbb{R} &\to \mathbb{R} \\
x &\mapsto x^2 .
\end{aligned}
\]
La imagen inversa de $y=1$ es $f^{-1}(1) = \{-1,1\}$. La imagen inversa de
$y=3$ es $f^{-1}(3) = \{-\sqrt{3},\sqrt{3}\}$. La imagen inversa de $y=-2$ es
$f^{-1}(-2) = \varnothing$.
\end{ejemplo}

\begin{ejemplo}
Sea el mapeo
\[
\begin{aligned}
f : \mathbb{R} &\to \mathbb{R} \\
x &\mapsto \sin x .
\end{aligned}
\]
La imagen inversa de $y=1$ es
\[
f^{-1}(1) = \left\{ (4n+1)\frac{\pi}{2} \; \middle| \; n \in \mathbb{Z} \right\}.
\]
La imagen inversa de $y=-1.5$ es $f^{-1}(-1.5) = \varnothing$.
\end{ejemplo}

\begin{definicion}[Dominio, rango e imagen]
El conjunto $X$ se denomina \textbf{dominio} del mapeo, mientras que $Y$ se
denomina \textbf{rango} del mapeo. Un mapeo no se puede definir sin
especificar el dominio y el rango.

La \textbf{imagen} del mapeo es
\[
f(X) = \{ y \in Y \mid y = f(x) \text{ para algún } x \in X \} \subset Y .
\]
La imagen $f(X)$ también se denota por $\operatorname{im} f$.
\end{definicion}

\begin{ejemplo}
\begin{enumerate}[label=\alph*)]
\item Sea el mapeo $f : \mathbb{R} \to \mathbb{R}$, $x \mapsto f(x) = e^x$. En
este ejemplo tanto el dominio como el rango son $\mathbb{R}$. La imagen del
mapeo es $\operatorname{im} f = f(\mathbb{R}) = \mathbb{R}^+ \subset
\mathbb{R}$. El valor $f(x) = -1$ no tiene imagen inversa, es decir
$f^{-1}(-1) = \varnothing$. Si cambiamos el dominio y el rango de
$\mathbb{R}$ al plano complejo $\mathbb{C}$, entonces
$f^{-1}(-1) = \{(2n+1)\pi i \mid n \in \mathbb{Z}\}$.

\item Sea el mapeo $f : \mathbb{R} \to \mathbb{R}$, $x \mapsto f(x) =
\sin(x)$. El dominio y el rango son $\mathbb{R}$. La imagen del mapeo es
$f(\mathbb{R}) = [-1,1]$. La imagen inversa de $0$ es
$f^{-1}(0) = \{ n\pi \mid n \in \mathbb{Z}\}$. Si cambiamos el dominio y el
rango a $\mathbb{C}$, entonces $f(\mathbb{C}) = \mathbb{C}$.
\end{enumerate}
\end{ejemplo}

\section{Inyectividad, sobreyectividad y biyectividad}

\begin{definicion}
Si un mapeo satisface alguna condición adicional, entonces recibe un nombre
especial:
\begin{enumerate}[label=\alph*)]
\item Un mapeo $f : X \to Y$ se dice \textbf{inyectivo} (o uno a uno) si
$x \neq x' $ implica $f(x) \neq f(x')$.
\item Un mapeo $f : X \to Y$ se dice \textbf{sobreyectivo} (o sobre) si para
cada $y \in Y$ existe al menos un elemento $x \in X$ tal que $f(x) = y$.
\item Un mapeo $f : X \to Y$ se dice \textbf{biyectivo} si es tanto inyectivo
como sobreyectivo.
\end{enumerate}
\end{definicion}

\begin{ejemplo}
Un mapeo definido por $f : x \mapsto ax$, con $a \in \mathbb{R} - \{0\}$, es
biyectivo. El mapeo definido por $f : x \mapsto x^2$ no es inyectivo ni
sobreyectivo. El mapeo dado por $f : x \mapsto e^x$ es inyectivo pero no
sobreyectivo.
\end{ejemplo}

\section{Mapeo constante, restricción y composición}

\begin{definicion}
Un \textbf{mapeo constante} $c : X \to Y$ está definido por $c(x) = y_0$,
donde $y_0$ es un elemento fijo de $Y$ y $x$ es un elemento arbitrario de
$X$.

Dado un mapeo $f : X \to Y$, podemos considerar su \textbf{restricción} a
$A \subset X$, denotada por
\[
f|_A : A \to Y .
\]

Dados dos mapeos $f : X \to Y$ y $g : Y \to Z$, el \textbf{mapeo compuesto}
de $f$ y $g$ es el mapa $g \circ f : X \to Z$ definido por
$g \circ f (x) = g(f(x))$.
\end{definicion}

Diagramáticamente, la composición de dos funciones se representa como
\[
X \xrightarrow{\ f\ } Y \xrightarrow{\ g\ } Z, \qquad
x \mapsto f(x) \mapsto g(f(x)).
\]

\section{Mapeo inclusión e identidad}

\begin{definicion}
Si $A \subset X$, un \textbf{mapeo inclusión} $i : A \to X$ está definido
por $i(a) = a$ para cualquier $a \in A$. Un mapeo inclusión a menudo se
escribe como $i : A \hookrightarrow X$.

El \textbf{mapeo identidad} $\operatorname{id}_X : X \to X$ es un caso
especial de mapeo inclusión, para el que $A = X$.
\end{definicion}

\section{Mapeo inverso}

\begin{proposicionc}
Si $f : X \to Y$, definido por $f : x \mapsto f(x)$, es biyectivo, entonces
existe un mapeo inverso $f^{-1} : Y \to X$ tal que $f^{-1} : f(x) \mapsto x$,
que también es biyectivo. Los mapeos $f$ y $f^{-1}$ satisfacen
\[
f \circ f^{-1} = \operatorname{id}_Y, \qquad f^{-1} \circ f =
\operatorname{id}_X .
\]
\end{proposicionc}

\begin{observacion}
Si graficamos los mapeos $f$ y $f^{-1}$, estos son simétricos con respecto a
la recta $y=x$, es decir, los puntos de $f$ y $f^{-1}$ son equidistantes con
respecto al mapeo identidad $y=x$. Si $(a, b=f(a))$ es un punto en la gráfica
de $f$, entonces el punto $(b=f(a), a)$ debe ser un punto en la gráfica de la
función inversa $f^{-1}$.
\end{observacion}

\begin{center}
\begin{tikzpicture}[scale=0.9]
  \draw[-Stealth] (-0.3,0) -- (4.3,0) node[right] {$x$};
  \draw[-Stealth] (0,-0.3) -- (0,4.3) node[above] {$y$};
  \draw[dashed] (0,0) -- (4,4) node[above right] {\footnotesize $y=x$};
  \draw[thick,colordef] plot[domain=0.3:4,samples=40] (\x, {0.5*\x*\x/4+0.4});
  \node[colordef] at (3.6,3.9) {\footnotesize$f$};
  \draw[thick,colorej] plot[domain=0.4:4.2,samples=40] ({0.5*\x*\x/4+0.4},\x);
  \node[colorej] at (3.9,3.3) {\footnotesize$f^{-1}$};
\end{tikzpicture}
\end{center}

Existe una definición alternativa de inversa:

\begin{definicion}
Si $f : X \to Y$ y $g : Y \to X$ satisfacen $f \circ g = \operatorname{id}_Y$
y $g \circ f = \operatorname{id}_X$, entonces $f$ y $g$ son biyecciones. El
mapeo $g$ se dice que es la \textbf{inversa derecha} de $f$, y el mapeo $f$
se dice que es la \textbf{inversa izquierda} de $g$.
\end{definicion}

\begin{proposicionc}
Si $f : X \to Y$ y $g : Y \to X$ satisfacen $g \circ f = \operatorname{id}_X$,
entonces $f$ es inyectiva y $g$ es sobreyectiva. Si además se cumple
$f \circ g = \operatorname{id}_Y$, se obtiene que $f$ y $g$ son ambas
biyecciones.
\end{proposicionc}

\chapter{Producto Cartesiano}

\section{Definición y ejemplos}

El producto cartesiano entre dos conjuntos es una operación que produce un
nuevo conjunto cuyos elementos son todos los pares ordenados que pueden
formarse, tal que el primer elemento pertenece al primer conjunto y el
segundo elemento pertenece al segundo conjunto.

\begin{definicion}[Par ordenado]
Un par ordenado se compone de dos elementos $x$ e $y$, y se escribe
$(x,y)$, donde $x$ es el primer elemento e $y$ es el segundo elemento. Se
concluye que dos pares ordenados $(x,y)$ y $(z,w)$ son iguales si $x=z$ e
$y=w$.
\end{definicion}

\begin{definicion}[Producto cartesiano]
El \textbf{producto cartesiano} de dos conjuntos $A$ y $B$ es el conjunto,
denotado por $A \times B$, de todos los pares ordenados de la forma
\[
A \times B = \{ (a,b) \mid a \in A \wedge b \in B \}.
\]
\end{definicion}

\begin{ejemplo}
Sean los conjuntos $A = \{1,2,3,4\}$ y $B = \{2,4\}$. El producto cartesiano
$A \times B$ se obtiene colocando al conjunto $A$ como la abscisa y al
conjunto $B$ como la ordenada de un sistema coordenado bidimensional.

\begin{center}
\begin{tikzpicture}[scale=0.8]
  \foreach \x in {1,2,3,4} \draw (\x,0.1) -- (\x,-0.1) node[below] {$\x$};
  \foreach \y in {2,4} \draw (0.1,\y) -- (-0.1,\y) node[left] {$\y$};
  \draw[-Stealth] (0,0) -- (5,0) node[right] {\footnotesize$A$};
  \draw[-Stealth] (0,0) -- (0,5) node[above] {\footnotesize$B$};
  \foreach \x in {1,2,3,4}{
    \foreach \y in {2,4}{
      \fill[colordef] (\x,\y) circle (2pt);
    }
  }
  \node at (4.6,4.6) {\footnotesize $A\times B$};
\end{tikzpicture}
\end{center}

lo que resulta en
\[
A \times B = \{(1,2),(1,4),(2,2),(2,4),(3,2),(3,4),(4,2),(4,4)\}.
\]

Ahora efectuamos el producto cartesiano $B \times A$, colocando esta vez a
$B$ como abscisa y a $A$ como ordenada:
\[
B \times A = \{(2,1),(2,2),(2,3),(2,4),(4,1),(4,2),(4,3),(4,4)\}.
\]
Es claro que $A \times B \neq B \times A$, por lo tanto el producto
cartesiano, en general, \textbf{no conmuta}.
\end{ejemplo}

\section{Cardinalidad del producto cartesiano}

Una pregunta natural es: ¿cuántos elementos posee el conjunto $A \times B$?
En otras palabras, ¿cuál es la cardinalidad de $A \times B$? Sea
$\#(A\times B)$ el número de elementos que contiene $A \times B$.

Sea $A$ un conjunto con $N_A$ elementos, $\#(A) = N_A$, y sea $B$ un
conjunto con $N_B$ elementos, $\#(B) = N_B$. Para encontrar la cardinalidad
de $A \times B$ contamos los elementos que contiene, asumiendo que cada
elemento aparece una sola vez. Tomando el primer elemento de $A$ y formando
todos los pares ordenados posibles con los elementos de $B$, se obtienen
$N_B$ pares ordenados. Tomando el segundo elemento de $A$ se obtienen otros
$N_B$ pares ordenados, y así sucesivamente hasta agotar los $N_A$ elementos
de $A$. Por lo tanto,
\[
\#(A \times B) = \underbrace{N_B + N_B + \dots + N_B}_{N_A} = N_A \cdot N_B .
\]

\begin{proposicionc}
\[
\#(A \times B) = N_A \cdot N_B = \#(A)\cdot\#(B).
\]
\end{proposicionc}

\begin{ejemplo}
Para los conjuntos $A=\{1,2,3,4\}$ y $B=\{2,4\}$ se tiene $\#(A) = 4$ y
$\#(B) = 2$, luego $\#(A\times B) = 4\cdot 2 = 8$.
\end{ejemplo}

\begin{proposicionc}
En general $A \times B \neq B \times A$. Sin embargo, ambos conjuntos
poseen el mismo número de elementos:
\[
\#(A \times B) = N_A \cdot N_B = N_B \cdot N_A = \#(B \times A).
\]
\end{proposicionc}

\subsection{Casos particulares}

Un caso extremo aparece cuando uno de los conjuntos es el conjunto vacío,
$\varnothing$, con $\#(\varnothing) = 0$, por lo que
\[
\#(A \times \varnothing) = N_A \cdot 0 = 0.
\]
El conjunto vacío actúa como el ``cero'' del producto cartesiano. Sin
embargo, si $A$ es un conjunto continuo, entonces $\#(A) = \infty$, y
\[
\#(A \times \varnothing) = \infty \cdot 0 = \; ?
\]
el producto no está determinado por conteo directo. Geométricamente, el
conjunto $A \times \varnothing$ se visualiza como un rectángulo de base $A$
y altura $\varnothing$, es decir, una altura sin elementos; por ello es
razonable que el rectángulo no contenga elementos.

\begin{definicion}
El producto cartesiano de un conjunto $A$ con el conjunto vacío es vacío:
\[
A \times \varnothing = \varnothing .
\]
\end{definicion}

Si uno de los conjuntos posee un solo elemento, por ejemplo $\#(A) = 1$,
entonces
\[
\#(A\times B) = 1\cdot N_B = N_B = \#(B),
\]
por lo tanto, un conjunto con un solo elemento actúa como la \textbf{identidad}
del producto cartesiano.

\begin{axiomac}
Sean $A$ y $B$ dos conjuntos, entonces existe un único conjunto
$A \times B$.
\end{axiomac}

\begin{proposicionc}
Sean $A \neq \varnothing$ y $B \neq \varnothing$, tales que
$A \neq B \Rightarrow A \times B \neq B \times A$.
\end{proposicionc}

\begin{proposicionc}
Sean $A \neq \varnothing$ y $B \neq \varnothing$. Si $A = B \Leftrightarrow
A \times B = B \times A$.
\end{proposicionc}

\begin{proposicionc}
Si $A$, $B$ y $C$ son conjuntos, entonces
\[
(A - B) \times C = (A \times C) - (B \times C).
\]
\end{proposicionc}

\begin{proposicionc}
Si $A \subseteq A'$ y $B \subseteq B'$, entonces $A \times B \subseteq A'
\times B'$.
\end{proposicionc}

\section{Producto cartesiano de conjuntos continuos}

Los casos interesantes aparecen cuando los conjuntos son continuos, de modo
que $\#(A \times B) = \infty$. El producto cartesiano entre conjuntos
continuos genera espacios geométricos interesantes.

\begin{ejemplo}
Sean los intervalos cerrados $A = [1,3]$ y $B = [2,4]$, subconjuntos de
$\mathbb{R}$. El producto cartesiano $A \times B$ es
\[
A \times B = \{ (x,y) \in \mathbb{R}^2 \mid 1 \le x \le 3 \wedge 2 \le y \le 4 \},
\]
que corresponde a un cuadrado compacto en $\mathbb{R}^2$.

\begin{center}
\begin{tikzpicture}[scale=0.9]
  \draw[-Stealth] (0,0) -- (5,0) node[right] {$x$};
  \draw[-Stealth] (0,0) -- (0,5) node[above] {$y$};
  \draw[fill=colordef!25,draw=colordef] (1,2) rectangle (3,4);
  \foreach \v/\l in {1/1,3/3} \draw (\v,0.08) -- (\v,-0.08) node[below] {$\l$};
  \foreach \v/\l in {2/2,4/4} \draw (0.08,\v) -- (-0.08,\v) node[left] {$\l$};
  \node at (2,3) {\footnotesize$A \times B$};
\end{tikzpicture}
\end{center}

Si en cambio se toma $A=[1,3[$ (semiabierto), el conjunto $A$ no contiene a
$x=3$, y esto se refleja en que $A \times B$ no es compacto: el lado
$\{3\} \times [2,4]$ queda excluido de $A \times B$.
\end{ejemplo}

\begin{ejemplo}
Consideremos el producto cartesiano entre un intervalo $A = [a,b] \subset
\mathbb{R}$ y un punto $B = \{c\} \subset \mathbb{R}$:
\[
A \times B = \{(x,y)\in\mathbb{R}^2 \mid a \le x \le b \wedge y = c\}
= \{(x,c)\in \mathbb{R}^2 \mid a\le x\le b\},
\]
que es un segmento de recta horizontal a la altura $y=c$. Análogamente,
$B\times A = \{(c,y)\in\mathbb{R}^2 \mid a\le y \le b\}$ es un segmento
vertical. En general, el producto cartesiano entre un segmento de
$\mathbb{R}$ y un punto genera un segmento de recta en $\mathbb{R}^2$.
\end{ejemplo}

\begin{observacion}
El producto cartesiano entre dos subconjuntos de $\mathbb{R}$ siempre genera
una región rectangular en $\mathbb{R}^2$; por lo tanto, una región de
$\mathbb{R}^2$ que no sea un rectángulo (por ejemplo, el disco unitario
$D = \{(x,y)\in\mathbb{R}^2 \mid x^2+y^2 \le 1\}$) no puede escribirse en la
forma $A \times B$.
\end{observacion}

\section{Producto cartesiano de conjuntos unidimensionales curvos}

El producto cartesiano entre subconjuntos unidimensionales curvos también
da origen a conjuntos en dimensiones mayores. Sea la circunferencia unitaria
\[
S^1 = \{(x,y)\in\mathbb{R}^2 \mid x^2+y^2=1\}
\]
y sea $I = [0,1]$ (usaremos la notación $I$ para un intervalo cerrado de
$\mathbb{R}$ de aquí en adelante).

\begin{ejemplo}
Para encontrar $S^1 \times I$ colocamos $S^1$ como base, horizontalmente en
el plano $Oxy$ de un sistema tridimensional, y en cada punto de $S^1$
colocamos el segmento $I$ en la dirección $z$:
\[
S^1 \times I = \{(x,y,z)\in\mathbb{R}^3 \mid x^2+y^2=1 \wedge 0\le z\le 1\}.
\]
El resultado es el manto de un cilindro, sin las tapas.

\begin{center}
\begin{tikzpicture}[scale=1]
  \draw[colordef, thick] (0,0) ellipse (1.4 and 0.5);
  \draw[colordef, thick] (-1.4,0) -- (-1.4,2.2);
  \draw[colordef, thick] (1.4,0) -- (1.4,2.2);
  \draw[colordef, thick] (0,2.2) ellipse (1.4 and 0.5);
  \node at (2.2,1.1) {\footnotesize$S^1\times I$};
\end{tikzpicture}
\end{center}

También se puede calcular $I \times S^1$, obteniendo
\[
I \times S^1 = \{(x,y,z)\in\mathbb{R}^3 \mid 0\le x\le 1 \wedge y^2+z^2=1\},
\]
nuevamente el manto de un cilindro, pero rotado $90^\circ$ respecto a
$S^1\times I$.
\end{ejemplo}

\begin{ejemplo}[El toro]
El producto cartesiano entre dos circunferencias $S^1$ genera el llamado
\textbf{toro}, denotado
\[
T^2 = S^1 \times S^1 .
\]
Se construye colocando una circunferencia $S^1$ (por ejemplo horizontal) y,
en cada punto de ella, otra circunferencia $S^1$, generando la superficie
de una cámara de bicicleta. Equivalentemente, se puede colocar primero
$S^1$ en forma vertical y en cada punto ubicar otra circunferencia $S^1$;
el resultado es también un toro.

\begin{center}
\begin{tikzpicture}[scale=1.1]
  \draw[colordef, thick] (0,0) ellipse (2.2 and 1.1);
  \draw[colordef, thick] (0,0) ellipse (1.1 and 0.55);
  \draw[colordef, thick] (-2.2,0) .. controls (-2.2,0.55) and (-1.1,0.55) .. (-1.1,0);
  \draw[colordef, thick, dashed] (-2.2,0) .. controls (-2.2,-0.55) and (-1.1,-0.55) .. (-1.1,0);
  \draw[colordef, thick] (2.2,0) .. controls (2.2,0.55) and (1.1,0.55) .. (1.1,0);
  \draw[colordef, thick, dashed] (2.2,0) .. controls (2.2,-0.55) and (1.1,-0.55) .. (1.1,0);
  \node at (0,-1.5) {\footnotesize$T^2 = S^1\times S^1$};
\end{tikzpicture}
\end{center}
\end{ejemplo}

\section{Producto cartesiano en dimensiones mayores}

El producto cartesiano $A\times B$ estudiado hasta ahora se ha realizado con
conjuntos unidimensionales, generando conjuntos bidimensionales. Sin
embargo, también se pueden realizar productos cartesianos entre conjuntos de
mayor dimensión.

\begin{ejemplo}
Sea $A \subset \mathbb{R}^2$ y sea $I=[a,b]\subset\mathbb{R}$. El producto
cartesiano $A \times I$ es
\[
A \times I = \{(x,y,z)\in\mathbb{R}^3 \mid (x,y)\in A \wedge a\le z\le b\},
\]
que genera un cilindro macizo en $\mathbb{R}^3$, es decir, un volumen.
\end{ejemplo}

\begin{observacion}
En el producto $A \times B$, si uno (o ambos) de los conjuntos no es
compacto, entonces $A\times B$ tampoco es compacto.
\end{observacion}

\section{Integración y producto cartesiano}

Dada una función $F(x,y)$ definida en $D \subset \mathbb{R}^2$, el cálculo de
la integral de $F$ sobre $D$ está dado por
\[
I = \int_D F(x,y)\, d^2x, \qquad d^2x = dx\,dy,
\]
donde $D$ se denomina \textbf{región de integración}.

Una región general $D$ se puede escribir en la forma
\[
D := \{(x,y)\in\mathbb{R}^2 \mid x_i \le x \le x_f \wedge f_1(x) \le y \le
f_2(x)\},
\]
de modo que la integral doble se escribe
\[
I = \int_D F(x,y)\, d^2x = \int_{x_i}^{x_f} \int_{f_1(x)}^{f_2(x)}
F(x,y)\, dy\, dx .
\]
La integral se evalúa en dos iteraciones: primero se evalúa la integral
interna (en $y$) y luego la integral externa (en $x$). El proceso no es en
general conmutativo:
\[
\int_{x_i}^{x_f}\int_{f_1(x)}^{f_2(x)} F(x,y)\,dy\,dx \neq
\int_{f_1(x)}^{f_2(x)}\int_{x_i}^{x_f} F(x,y)\,dx\,dy .
\]
Para que el proceso sea conmutativo debemos despejar $y$ en función de $x$ y
reescribir la región $D$ como
\[
D := \{(x,y)\in\mathbb{R}^2 \mid y_i\le y\le y_f \wedge g_1(y)\le x\le
g_2(y)\},
\]
de modo que
\[
I = \int_D F(x,y)\, d^2x = \int_{y_i}^{y_f}\int_{g_1(y)}^{g_2(y)} F(x,y)\,
dx\, dy .
\]
Ambas formas de la integral deben coincidir:
\[
\int_{x_i}^{x_f}\int_{f_1(x)}^{f_2(x)} F(x,y)\,dy\,dx
= \int_{y_i}^{y_f}\int_{g_1(y)}^{g_2(y)} F(x,y)\, dx\, dy .
\]

\begin{ejemplo}[Área del disco compacto]
Sea el disco compacto $D = \{(x,y)\in\mathbb{R}^2 \mid x^2+y^2\le r^2\}$.
Calculemos su área,
\[
A = \int_D dx\,dy .
\]
Reescribiendo $D$ en la forma
\[
D = \{(x,y)\in\mathbb{R}^2 \mid -r\le x\le r \wedge -\sqrt{r^2-x^2}\le y \le
\sqrt{r^2-x^2}\},
\]
la integral queda
\[
A = \int_{-r}^{r} \int_{-\sqrt{r^2-x^2}}^{\sqrt{r^2-x^2}} dy\,dx
= \int_{-r}^{r} y \Big|_{-\sqrt{r^2-x^2}}^{\sqrt{r^2-x^2}} dx
= 2\int_{-r}^{r} \sqrt{r^2-x^2}\, dx .
\]
Con el cambio de variable $x(\alpha) = r\sin\alpha$,
\[
A = 2r^2 \int_{-\pi/2}^{\pi/2} \cos^2\alpha\, d\alpha
= r^2 \left(\alpha + \tfrac12\sin(2\alpha)\right)\Big|_{-\pi/2}^{\pi/2}
\;\Longrightarrow\; A = \pi r^2 .
\]
\end{ejemplo}

\subsection{El caso rectangular \texorpdfstring{$D=I_1\times I_2$}{D=I1xI2}}

El caso de mayor interés es cuando $D = I_1\times I_2$ es un rectángulo en
$\mathbb{R}^2$,
\[
D = I_1\times I_2 = \{(x,y)\in\mathbb{R}^2 \mid x_i\le x\le x_f \wedge
y_i\le y\le y_f\},
\]
de modo que la integral $I$ toma la forma
\[
I = \int_D F(x,y)\, d^2x = \int_{I_1\times I_2} F(x,y)\, d^2x
= \int_{I_1}\int_{I_2} F(x,y)\, dy\, dx
= \int_{x_i}^{x_f}\int_{y_i}^{y_f} F(x,y)\,dy\,dx .
\]
En este caso particular sí se cumple
\[
\int_{I_1\times I_2} F(x,y)\, d^2x = \int_{I_1}\int_{I_2} F(x,y)\,dy\,dx
= \int_{I_2}\int_{I_1} F(x,y)\,dy\,dx = \int_{I_2\times I_1} F(x,y)\, d^2x ,
\]
aun cuando, en general, $I_1\times I_2 \neq I_2\times I_1$ como conjuntos
ordenados.

\begin{observacion}
La integración sobre una región $D=I_1\times I_2$ es distinta de una
integración sobre una región $D = D_1 \cup D_2$, puesto que
\[
I = \int_{I_1\times I_2} F(x,y)\, d^2x = \int_{I_1}\int_{I_2} F(x,y)\,dy\,dx,
\]
\[
I' = \int_{D_1\cup D_2} F(x,y)\, d^2x = \int_{D_1} F(x,y)\,d^2x +
\int_{D_2} F(x,y)\, d^2x,
\]
donde $I_i\subset\mathbb{R}$ y $D_i\subset\mathbb{R}^2$, con $i=1,2$.
\end{observacion}

\subsection{Integrales de una derivada parcial y el borde del rectángulo}

Consideremos el caso $F(x,y) = \dfrac{\partial g(x,y)}{\partial y}$. Entonces
\[
I = \int_D F(x,y)\, d^2x = \int_{I_1\times I_2} F(x,y)\,dy\,dx
= \int_{x_i}^{x_f}\int_{y_i}^{y_f} \frac{\partial g(x,y)}{\partial y}\,
dy\, dx
= \int_{x_i}^{x_f} g(x,y)\Big|_{y=y_i}^{y=y_f} dx,
\]
es decir,
\[
I = \int_{x_i}^{x_f} g(x,y_f)\, dx - \int_{x_i}^{x_f} g(x,y_i)\, dx .
\]

La primera integral está realizada sobre la región $I_1\times\{y_f\}$, que
es uno de los bordes del rectángulo $M = I_1\times I_2$. La segunda integral
está realizada sobre $I_1\times\{y_i\}$, el borde opuesto.

\begin{center}
\begin{tikzpicture}[scale=0.9]
  \draw[-Stealth] (0,0) -- (5,0) node[right] {$x$};
  \draw[-Stealth] (0,0) -- (0,4) node[above] {$y$};
  \draw[fill=colordef!12,draw=colordef] (0.8,0.8) rectangle (3.8,2.8);
  \draw[colorej, ultra thick] (0.8,2.8) -- (3.8,2.8);
  \draw[colorej, ultra thick] (0.8,0.8) -- (3.8,0.8);
  \node at (2.3,3.3) {\footnotesize$I_1\times\{y_f\}$};
  \node at (2.3,0.3) {\footnotesize$I_1\times\{y_i\}$};
  \node at (2.3,1.8) {\footnotesize$M$};
\end{tikzpicture}
\end{center}

Ambas integrales se pueden reescribir en la forma
\[
I = \int_{x_i}^{x_f} g(x,y_f)\,dx - \int_{x_i}^{x_f} g(x,y_i)\,dx
= \int_{I_1} g(x,y_f)\,dx - \int_{I_1} g(x,y_i)\,dx
= \int_{I_1\times\{y_f\}} g(x,y)\,dx - \int_{I_1\times\{y_i\}} g(x,y)\,dx,
\]
donde los puntos $\{y_i\}$ e $\{y_f\}$ actúan aquí como una ``integral
interior''. Esta observación anticipa la relación entre integrales de
derivadas y la frontera de la región de integración, que se desarrolla en el
capítulo siguiente.

\chapter{Frontera}

\section{Espacios topológicos}

\begin{definicion}[Espacio topológico]
Sea $X$ un conjunto y sea $T = \{U_i \mid i \in I\}$, donde $I$ es un
conjunto de índices, una colección de subconjuntos de $X$, es decir
$U_i \subset X$. El par $(X,T)$ se dice un \textbf{espacio topológico} si
$T$ satisface los siguientes axiomas:
\begin{enumerate}[label=\roman*)]
\item $\varnothing, X \in T$.
\item La unión arbitraria de elementos de $T$ es un elemento de $T$:
\[
\bigcup_{i\in I} U_i \in T .
\]
\item La intersección finita de elementos de $T$ es un elemento de $T$:
\[
\bigcap_{\text{finita}} U_i \in T .
\]
\end{enumerate}
Los elementos $U_i$ de $T$ se llaman \textbf{conjuntos abiertos}.
\end{definicion}

\begin{observacion}
El conjunto $T$ es un conjunto de conjuntos: los elementos de $T$ son, a su
vez, conjuntos (los subconjuntos abiertos de $X$).
\end{observacion}

\begin{ejemplo}
Sea $X = \{a,b,c,d,e\}$.
\begin{itemize}
\item $T_1 = \{\varnothing, X, \{a\}, \{c,d\}, \{a,c,d\}, \{b,c,d,e\}\}$
\textbf{es} una topología sobre $X$.
\item $T_2 = \{\varnothing, X, \{a\}, \{c,d\}, \{a,c,d\}, \{a,b,c,d\}\}$
\textbf{es} una topología sobre $X$.
\item $T_3 = \{\varnothing, X, \{a\}, \{e\}, \{a,e\}, \{a,c,e\},
\{b,c,d,e\}\}$ \textbf{no es} una topología sobre $X$, puesto que
\[
\{a,c,e\} \cap \{b,c,d,e\} = \{c,e\} \notin T_3 .
\]
\end{itemize}
\end{ejemplo}

\begin{ejemplo}
Sea $X = \mathbb{R}$.
\begin{itemize}
\item $T_1$, compuesta por $\mathbb{R}$, $\varnothing$ y todos los
intervalos $(-n,n)$ con $n\in\mathbb{N}$, es una topología sobre
$\mathbb{R}$.
\item $T_2$, compuesta por $\mathbb{R}$, $\varnothing$ y todos los
intervalos $[-n,n]$ con $n\in\mathbb{N}$, es una topología sobre
$\mathbb{R}$.
\item $T_3$, compuesta por $\mathbb{R}$, $\varnothing$ y todos los
intervalos $[n,\infty)$ con $n\in\mathbb{N}$, es una topología sobre
$\mathbb{R}$.
\item $T_4$, compuesta por $\mathbb{R}$, $\varnothing$ y todos los
intervalos $(a,b)$, $\forall a,b\in\mathbb{R}$ con $a<b$, junto con todas sus
uniones arbitrarias e intersecciones finitas, es una topología sobre
$\mathbb{R}$ llamada \textbf{topología usual} o \textbf{canónica}.
\end{itemize}
\end{ejemplo}

\begin{definicion}[Topología discreta y trivial]
Si $X$ es un conjunto y $T$ es la colección de \emph{todos} los
subconjuntos de $X$, entonces $T$ es una topología de $X$ llamada
\textbf{topología discreta}. Si $T = \{\varnothing, X\}$, entonces $T$ es
una topología de $X$ llamada \textbf{topología trivial}.
\end{definicion}

\section{Espacios de Hausdorff}

\begin{definicion}[Separabilidad / Hausdorff]
Un espacio topológico $(X,T)$ se dice \textbf{separable} o de
\textbf{Hausdorff} si para $x_1,x_2\in X$, con $x_1\neq x_2$, existen dos
abiertos $x_1\in U_1$ y $x_2\in U_2$ tales que
\[
U_1 \cap U_2 = \varnothing .
\]
\end{definicion}

\begin{ejemplo}
Sea $X=\mathbb{R}^2$ y consideremos dos topologías:
\begin{enumerate}[label=\alph*)]
\item $T_1$, compuesta por $\mathbb{R}^2$, $\varnothing$ y todos los
conjuntos
\[
D_r := \{(x,y)\in\mathbb{R}^2 \mid x^2+y^2 < r^2\}, \qquad r>0,
\]
es una topología de $X=\mathbb{R}^2$ (discos centrados en el origen).

\item $T_2$, compuesta por $\mathbb{R}^2$, $\varnothing$ y todos los
conjuntos
\[
D_r(a,b) := \{(x,y)\in\mathbb{R}^2 \mid (x-a)^2+(y-b)^2 < r^2\}, \qquad r>0,\
(a,b)\in\mathbb{R}^2,
\]
junto con todas sus uniones arbitrarias e intersecciones finitas, es la
\textbf{topología usual} de $X=\mathbb{R}^2$ (discos centrados en cualquier
punto).
\end{enumerate}

\begin{center}
\begin{tikzpicture}[scale=0.85]
  \begin{scope}
    \draw[-Stealth] (0,0) -- (3.3,0) node[right] {$x$};
    \draw[-Stealth] (0,0) -- (0,3.3) node[above] {$y$};
    \draw[colordef] (0,0) circle (1.4);
    \fill (0,0) circle (1pt);
    \node at (0,-0.35) {\footnotesize$(0,0)$};
    \node at (1.6,-0.4) {\footnotesize$T_1$: discos en el origen};
  \end{scope}
  \begin{scope}[xshift=6.5cm]
    \draw[-Stealth] (0,0) -- (3.3,0) node[right] {$x$};
    \draw[-Stealth] (0,0) -- (0,3.3) node[above] {$y$};
    \draw[colorej] (1,1.7) circle (0.9);
    \fill (1,1.7) circle (1pt) node[above right] {\footnotesize$(a,b)$};
    \node at (1.6,-0.4) {\footnotesize$T_2$: discos en $(a,b)$};
  \end{scope}
\end{tikzpicture}
\end{center}

Veamos cuál de los dos espacios topológicos, $(\mathbb{R}^2,T_1)$ o
$(\mathbb{R}^2,T_2)$, es de Hausdorff.

\medskip
\noindent\textbf{Caso $(\mathbb{R}^2,T_1)$:} sean $x_1=(1,0)$,
$x_2=(0,2)\in\mathbb{R}^2$, y escojamos los abiertos
\[
D_{1.1} = \{(x,y)\in\mathbb{R}^2 \mid x^2+y^2 < (1.1)^2\}, \qquad
D_{2.1} = \{(x,y)\in\mathbb{R}^2 \mid x^2+y^2 < (2.1)^2\},
\]
con $x_1\in D_{1.1}$ y $x_2\in D_{2.1}$. Como $D_{1.1}\subset D_{2.1}$,
se tiene $D_{1.1}\cap D_{2.1} = D_{1.1} \neq \varnothing$. Como cualquier
par de abiertos centrados en el origen que contengan a $x_1$ y $x_2$ se
intersectan de la misma forma, se concluye que $(\mathbb{R}^2,T_1)$
\textbf{no} es un espacio de Hausdorff.

\medskip
\noindent\textbf{Caso $(\mathbb{R}^2,T_2)$:} con los mismos puntos,
escojamos
\[
D_1(1,0) = \{(x,y)\in\mathbb{R}^2 \mid (x-1)^2+y^2<1\}, \qquad
D_1(0,2) = \{(x,y)\in\mathbb{R}^2 \mid x^2+(y-2)^2<1\},
\]
con $x_1\in D_1(1,0)$ y $x_2\in D_1(0,2)$. Se verifica que
$D_1(1,0)\cap D_1(0,2) = \varnothing$, por lo tanto $(\mathbb{R}^2,T_2)$
\textbf{sí} es un espacio de Hausdorff.
\end{ejemplo}

\section{Entornos}

Antes de continuar necesitaremos una definición elemental de teoría de
conjuntos.

\begin{definicion}
Sean los conjuntos $A$ y $B$. Se define el conjunto $A-B$ como
\[
A - B = \{ x \mid x\in A \wedge x\notin B \}.
\]
\end{definicion}

\begin{definicion}[Entorno]
Sea $X$ un espacio topológico y sea $S\subseteq X$. Si $x\in S$, se dice
que $V_x \subset X$ es un \textbf{entorno} del punto $x$ si contiene un
conjunto abierto $U_x$ que a su vez contiene a $x$:
\[
x \in U_x \subseteq V_x .
\]
\end{definicion}

\begin{ejemplo}
Sea $S=[0,1]\subset\mathbb{R}$ con la topología usual. Entonces
$[0,1)$ es un entorno de $x=1/2$, pues contiene al abierto $(1/4,3/4)$.
También $(1/4,3/4)$ es en sí mismo un entorno de $x=1/2$ (y además es
abierto).

Sea $X=\mathbb{R}^2$ con la topología usual. Un entorno del punto
$x=(0,1)$ es
\[
V_x = \{(x,y)\in\mathbb{R}^2 \mid -1\le x\le 1 \wedge -2\le y\le 2\},
\]
puesto que contiene al abierto
\[
U_x = \{(x,y)\in\mathbb{R}^2 \mid x^2+(y-1)^2 < 1\}.
\]
\end{ejemplo}

\begin{observacion}
El entorno $V_x$ no necesariamente es un conjunto abierto. Si $V_x$ es
abierto se le llama \textbf{entorno abierto}. Algunos autores exigen que
los entornos sean siempre abiertos, por lo que conviene prestar atención a
la convención de cada texto.
\end{observacion}

\section{Un teorema auxiliar: la propiedad de Arquímedes}

\begin{proposicionc}[Arquímedes]
Sean dos puntos $x,y\in\mathbb{R}$ con $x<y$. Siempre existe al menos un
punto $z\in\mathbb{R}$ tal que $x<z<y$.
\end{proposicionc}

Damos una demostración artesanal, no rigurosa pero ilustrativa. Sean $x<y$;
un punto entre ambos es el punto medio $z = \tfrac{x+y}{2}$. En efecto,
\[
x<y \;\big/+x \;\Rightarrow\; 2x<x+y \;\Rightarrow\; x<\frac{x+y}{2},
\]
\[
x<y \;\big/+y \;\Rightarrow\; x+y<2y \;\Rightarrow\; \frac{x+y}{2}<y,
\]
por lo tanto
\[
x < \frac{x+y}{2} < y .
\]
Este argumento se puede iterar: el punto medio entre $x$ y $\tfrac{x+y}{2}$
es
\[
\frac{x+\tfrac{x+y}{2}}{2} = \frac{3x+y}{4}, \qquad\text{y se cumple}\qquad
x < \frac{3x+y}{4} < \frac{x+y}{2},
\]
generando una sucesión de puntos intermedios cada vez más cercanos a $x$.

\begin{definicion}
Se define $\min\{x_1,x_2,\dots,x_n\}$ como la menor de las cantidades
$x_1,x_2,\dots,x_n \in \mathbb{R}$.
\end{definicion}

\begin{ejemplo}
$\min\{1,3\} = 1$. \qquad $\min\{1,-3\} = -3$.
\end{ejemplo}

\section{Puntos de acumulación}

\begin{definicion}[Punto de acumulación]
Sea $(X,T)$ un espacio topológico y $S$ un subconjunto de $X$. Diremos que
$x\in X$ es un \textbf{punto de acumulación} (o \textbf{punto límite}) de
$S$ si y solo si para cualquier abierto $U_x$ que sea entorno de $x$ se
cumple
\[
S \cap (U_x - \{x\}) \neq \varnothing .
\]
\end{definicion}

\begin{ejemplo}
Sea $S=[0,1]\subset\mathbb{R}$ con la topología usual.

\emph{Puntos interiores.} Sea $x\in S$ con $0<x<1$, y sea el entorno abierto
$U_x = (x-\epsilon,x+\epsilon)$ con
\[
0 < \epsilon < \min\left\{ \frac{x}{2}, \frac{1-x}{2}\right\}
\]
(por la propiedad de Arquímedes). Entonces
\[
S\cap(U_x-\{x\}) = [0,1]\cap\big((x-\epsilon,x)\cup(x,x+\epsilon)\big)
= (x-\epsilon,x)\cup(x,x+\epsilon) \neq \varnothing,
\]
por lo tanto, todos los puntos $0<x<1$ de $S$ son puntos de acumulación de
$S$.

\emph{El punto $x=0$.} Tomando el abierto $(-\epsilon,\epsilon)$ con
$\epsilon$ pequeño (por ejemplo $\epsilon = 10^{-5}$):
\[
S\cap(U_x-\{0\}) = [0,1]\cap\big((-\epsilon,0)\cup(0,\epsilon)\big) =
(0,\epsilon) \neq \varnothing,
\]
de modo que $x=0$ es punto de acumulación de $S$. Análogamente, $x=1$ es
punto de acumulación de $S$.

\begin{center}
\begin{tikzpicture}[scale=1.1]
  \draw[thick] (0,0) -- (4,0);
  \draw[fill=black] (0,0) circle (2pt) node[below] {\footnotesize$0$};
  \draw[fill=black] (2.6,0) circle (2pt) node[below] {\footnotesize$1$};
  \draw (-0.4,0) node[left] {\footnotesize$($};
  \draw[decorate, decoration={brace, amplitude=4pt}] (-0.15,0.25) -- (0.55,0.25);
  \node at (0.2,0.6) {\footnotesize$\epsilon$};
\end{tikzpicture}
\end{center}

\emph{Puntos fuera de $S$.} Sea $x=1+\epsilon \notin S$, con $\epsilon>0$
muy pequeño. Consideremos el abierto
\[
U_x = \left(1+\frac{\epsilon}{2},\; 1+\frac{3\epsilon}{2}\right),
\]
centrado en $x=1+\epsilon$. Se concluye que
\[
S\cap(U_x - \{1+\epsilon\}) = [0,1]\cap
\left(\left(1+\tfrac{\epsilon}{2},1+\epsilon\right)\cup
\left(1+\epsilon,1+\tfrac{3\epsilon}{2}\right)\right) = \varnothing,
\]
por lo que $x=1+\epsilon$ \textbf{no} es punto de acumulación de $S$. Lo
mismo ocurre para $x=-\epsilon<0$. En general, los puntos $x<0$ y $x>1$ no
son puntos de acumulación de $S$.
\end{ejemplo}

\begin{ejemplo}
Sea ahora $S=(0,1)\subset\mathbb{R}$ con la topología usual. Cada punto
interior $0<x<1$ es punto de acumulación (mismo argumento que antes). El
punto $x=0$ también es punto de acumulación, pues para el abierto
$(-\epsilon,\epsilon)$,
\[
S\cap(U_x-\{0\}) = (0,1)\cap\big((-\epsilon,0)\cup(0,\epsilon)\big) =
(0,\epsilon)\neq\varnothing.
\]
Lo mismo ocurre para $x=1$. Es decir, $x=0$ y $x=1$ son puntos de
acumulación de $S$ \textbf{aunque no pertenecen a $S$}. Los puntos $x<0$ y
$x>1$ no son puntos de acumulación.
\end{ejemplo}

\begin{observacion}
Los conjuntos $S_1=[0,1]$ y $S_2=(0,1)$ poseen los \emph{mismos} puntos de
acumulación, aun cuando $S_2 \subsetneq S_1$.
\end{observacion}

\begin{ejemplo}
Sea $S=\{(x,y)\in\mathbb{R}^2 \mid 0\le x\le 1 \wedge 0\le y\le 1\}
\subset\mathbb{R}^2$, con la topología usual de $\mathbb{R}^2$. El conjunto
$S$ se puede escribir como
\[
S = [0,1]\times[0,1] = I\times I = I^2,
\]
una región cuadrada en $\mathbb{R}^2$.

\begin{center}
\begin{tikzpicture}[scale=1.6]
  \draw[-Stealth] (0,0) -- (1.4,0) node[right] {$x$};
  \draw[-Stealth] (0,0) -- (0,1.4) node[above] {$y$};
  \draw[fill=colordef!15,draw=colordef] (0,0) rectangle (1,1);
  \node at (0,-0.13) {\footnotesize$0$};
  \node at (1,-0.13) {\footnotesize$1$};
  \node at (-0.13,1) {\footnotesize$1$};
\end{tikzpicture}
\end{center}

\emph{Puntos interiores.} Sea $x=(a,b)\in S$ con $0<a<1$ y $0<b<1$; es un
punto de acumulación de $S$. Basta tomar el entorno abierto
\[
U_x = \{(x,y)\in\mathbb{R}^2 \mid (x-a)^2+(y-b)^2<r^2\}, \qquad
r < \min\{1-a,\,a,\,1-b,\,b\},
\]
de modo que $U_x - \{(a,b)\} \subset S$, y por lo tanto
$S\cap(U_x-\{(a,b)\}) \neq \varnothing$.

\emph{Puntos del borde.} Sea $x=(a,0)$ con $0\le a\le 1$. Tomando el abierto
\[
U_x = \{(x,y)\in\mathbb{R}^2 \mid (x-a)^2+y^2<r^2\}, \qquad r<\min\{a,1-a\},
\]
se tiene $S\cap(U_x-\{(a,0)\}) \neq \varnothing$, de modo que todo el
segmento
\[
L_1 = \{(a,0)\in\mathbb{R}^2 \mid 0\le a\le 1\}
\]
está formado por puntos de acumulación de $S$. Del mismo modo son puntos
de acumulación de $S$ los segmentos
\[
L_2 = \{(1,b)\in\mathbb{R}^2 \mid 0\le b\le 1\}, \quad
L_3 = \{(a,1)\in\mathbb{R}^2 \mid 0\le a\le 1\}, \quad
L_4 = \{(0,b)\in\mathbb{R}^2 \mid 0\le b\le 1\}.
\]

\begin{center}
\begin{tikzpicture}[scale=1.6]
  \draw[-Stealth] (0,0) -- (1.4,0) node[right] {$x$};
  \draw[-Stealth] (0,0) -- (0,1.4) node[above] {$y$};
  \draw[fill=colordef!10,draw=none] (0,0) rectangle (1,1);
  \draw[colorej,ultra thick] (0,0) -- (1,0) node[midway,below] {\footnotesize$L_1$};
  \draw[colorej,ultra thick] (1,0) -- (1,1) node[midway,right] {\footnotesize$L_2$};
  \draw[colorej,ultra thick] (0,1) -- (1,1) node[midway,above] {\footnotesize$L_3$};
  \draw[colorej,ultra thick] (0,0) -- (0,1) node[midway,left] {\footnotesize$L_4$};
\end{tikzpicture}
\end{center}
\end{ejemplo}

\section{Interior de un conjunto}

\begin{definicion}[Interior]
Sea $S\subseteq X$. Se define el \textbf{conjunto interior} de $S$,
denotado $\operatorname{int}(S)$ o $\mathring{S}$, como la unión de todos
los abiertos contenidos en $S$.
\end{definicion}

\begin{ejemplo}
Sea $S=[0,1]\subset\mathbb{R}$ con la topología usual. El conjunto
$\{x\in\mathbb{R} \mid 0<x<1\}$ es el interior de $S$:
\[
\mathring{S} = \{x\in\mathbb{R}\mid 0<x<1\}.
\]
En efecto, sea $x$ con $0<x<1$ y sea el abierto
$U_x=(x-\epsilon,x+\epsilon)$ con $\epsilon<\min\{x,1-x\}$, tal que
$S\cap U_x = U_x$, luego $U_x\subset S$.

Sin embargo, el punto $x=0$ está contenido solo en abiertos que \emph{no}
están enteramente contenidos en $S$: tomando $U_x=(-\epsilon,\epsilon)$,
\[
S\cap U_x = (0,\epsilon),
\]
que si bien es abierto, no contiene al punto $x=0$; por lo tanto
$0\notin\mathring{S}$. Lo mismo ocurre con $x=1\in S$, así que
$1\notin\mathring{S}$.

Notar que para $S=(0,1)$ se tiene igualmente $\mathring{S}=(0,1)=S$.
\end{ejemplo}

\begin{proposicionc}
Sea $S\subseteq X$, entonces $\mathring{S}\subseteq S$.
\end{proposicionc}

\section{Conjuntos compactos}

\begin{definicion}[Compacidad]
Un conjunto $S\subseteq X$ se dice \textbf{compacto} si contiene a todos
sus puntos de acumulación.
\end{definicion}

\begin{ejemplo}
\begin{itemize}
\item $S=[0,1]$ es compacto, pues contiene a todos sus puntos de
acumulación.
\item $S=(0,1)$ \textbf{no} es compacto: no contiene a los puntos de
acumulación $x=0$ y $x=1$.
\item $S=[0,1)$ \textbf{no} es compacto: no contiene al punto de
acumulación $x=1$.
\item $S=\mathbb{R}$ \textbf{no} es compacto: no contiene a los puntos de
acumulación $x=-\infty$ y $x=+\infty$.
\end{itemize}
\end{ejemplo}

\section{Clausura de un conjunto}

\begin{definicion}[Clausura]
Sea $S\subseteq X$. Se llama \textbf{clausura} o \textbf{cierre} de $S$ al
conjunto de todos los puntos de acumulación de $S$, y se denota
$\overline{S}$.
\end{definicion}

\begin{ejemplo}
\begin{itemize}
\item $S=[0,1]\subset\mathbb{R}$: $\overline{S}=S$.
\item $S=(0,1)$: $\overline{S} = S\cup\{0\}\cup\{1\} = [0,1]$.
\item $S=[0,1)$: $\overline{S} = S\cup\{1\} = [0,1]$.
\item $S=[a,b]$, con $a<b$: $\overline{S}=S$.
\item $S=(a,b]$, con $a<b$: $\overline{S} = \{a\}\cup S$.
\item $S=\mathbb{R}$: $\overline{\mathbb{R}} = \mathbb{R}\cup(\{-\infty\}
\cup\{\infty\}) = [-\infty,\infty] = \{x \mid -\infty\le x\le\infty\}$.
\end{itemize}
\end{ejemplo}

\begin{proposicionc}
Sea $X$ un espacio topológico y $S\subset X$, entonces $S\subseteq
\overline{S}$.
\end{proposicionc}

\section{Frontera de un conjunto}

\begin{definicion}[Frontera]
Sea $X$ un espacio topológico y $S\subseteq X$. Se define el conjunto
\textbf{frontera} de $S$, denotado $\partial S$, como
\[
\partial S = \overline{S} \cap \overline{(X-S)} .
\]
La frontera también se denomina \textbf{borde}.
\end{definicion}

\begin{ejemplo}
Sea $S=[0,1]\subset\mathbb{R}$ con la topología usual. La clausura de $S$
es $\overline{S}=[0,1]=S$. El conjunto $\mathbb{R}-S$ es
\[
\mathbb{R}-S = (-\infty,0)\cup(1,+\infty),
\]
cuya clausura es
\[
\overline{(\mathbb{R}-S)} = (\mathbb{R}-S)\cup\{-\infty\}\cup\{0\}\cup
\{1\}\cup\{+\infty\} = [-\infty,0]\cup[1,+\infty].
\]
Entonces la frontera es
\[
\partial S = \overline{S}\cap\overline{(\mathbb{R}-S)} = [0,1]\cap
\big([-\infty,0]\cup[1,+\infty]\big) = \{0\}\cup\{1\}.
\]
\end{ejemplo}

\begin{ejemplo}
En general, sea $S=[a,b]$ con $a<b$. Análogamente al ejemplo anterior,
\[
\mathbb{R}-S = (-\infty,a)\cup(b,+\infty), \qquad
\overline{(\mathbb{R}-S)} = [-\infty,a]\cup[b,+\infty],
\]
por lo que
\[
\partial S = [a,b]\cap\big([-\infty,a]\cup[b,+\infty]\big) = \{a\}\cup\{b\}.
\]
\end{ejemplo}

\begin{ejemplo}
\begin{itemize}
\item $S=(0,1)$: $\partial S = \{0\}\cup\{1\}$.
\item $S=[0,1)$: $\partial S = \{0\}\cup\{1\}$.
\item $S=X=\mathbb{R}$: $\partial S = \{-\infty\}\cup\{+\infty\}$.
\end{itemize}
\end{ejemplo}

\begin{ejemplo}
Sea la circunferencia $S_1 = \{(x,y)\in\mathbb{R}^2\mid x^2+y^2=1\}$;
entonces $\partial S_1 = \varnothing$. Sea la esfera
$S_2 = \{(x,y,z)\in\mathbb{R}^3\mid x^2+y^2+z^2=1\}$; entonces
$\partial S_2 = \varnothing$. En general, para la $n$-esfera
\[
S^n = \{(x_1,\dots,x_{n+1})\in\mathbb{R}^{n+1} \mid x_1^2+\dots+x_{n+1}^2
=1\},
\]
se tiene $\partial S^n = \varnothing$.
\end{ejemplo}

\begin{definicion}[Conjunto sin frontera]
Un conjunto $S\subseteq X$, donde $X$ es un espacio topológico, se dice
\textbf{sin frontera} o \textbf{sin borde} si $\partial S = \varnothing$.
\end{definicion}

\begin{proposicionc}
Sea $X$ un espacio topológico continuo y sea $S\subseteq X$. Si $S$ es un
conjunto continuo (sin frontera), entonces
\[
\partial^2 S = \partial(\partial S) = \varnothing .
\]
\end{proposicionc}

\begin{proposicionc}
Sea $X$ un espacio topológico y sean $S_1,S_2\subseteq X$. Entonces
\[
\partial(S_1\cup S_2) = \partial S_1 \cup \partial S_2, \qquad
\partial(S_1\times S_2) = (\partial S_1 \times S_2) \cup
(S_1\times \partial S_2).
\]
\end{proposicionc}

\section{La frontera de un rectángulo}

\begin{ejemplo}
\label{ej:frontera-rectangulo}
Sea $S = I_1\times I_2 = [a,b]\times[c,d]$ con $a<b$ y $c<d$.

\begin{center}
\begin{tikzpicture}[scale=0.95]
  \draw[-Stealth] (0,0) -- (5,0) node[right] {$x$};
  \draw[-Stealth] (0,0) -- (0,4) node[above] {$y$};
  \draw[fill=colordef!15,draw=colordef] (1,0.8) rectangle (3.6,2.9);
  \node at (1,-0.3) {\footnotesize$a$};
  \node at (3.6,-0.3) {\footnotesize$b$};
  \node at (-0.3,0.8) {\footnotesize$c$};
  \node at (-0.3,2.9) {\footnotesize$d$};
\end{tikzpicture}
\end{center}

La frontera de $S$ es $\partial S = L_1\cup L_2\cup L_3\cup L_4$:

\begin{center}
\begin{tikzpicture}[scale=0.95]
  \draw[-Stealth] (0,0) -- (5,0) node[right] {$x$};
  \draw[-Stealth] (0,0) -- (0,4) node[above] {$y$};
  \draw[colordef!25,fill=colordef!8] (1,0.8) rectangle (3.6,2.9);
  \draw[colorej,ultra thick] (1,0.8) -- (3.6,0.8) node[midway,below] {\footnotesize$L_1$};
  \draw[colorej,ultra thick] (3.6,0.8) -- (3.6,2.9) node[midway,right] {\footnotesize$L_2$};
  \draw[colorej,ultra thick] (1,2.9) -- (3.6,2.9) node[midway,above] {\footnotesize$L_3$};
  \draw[colorej,ultra thick] (1,0.8) -- (1,2.9) node[midway,left] {\footnotesize$L_4$};
  \node at (1,-0.3) {\footnotesize$a$};
  \node at (3.6,-0.3) {\footnotesize$b$};
  \node at (-0.3,0.8) {\footnotesize$c$};
  \node at (-0.3,2.9) {\footnotesize$d$};
\end{tikzpicture}
\end{center}

donde
\[
\begin{aligned}
L_1 &= \{(x,c)\in\mathbb{R}^2 \mid a<x<b\}
= \{(x,y)\in\mathbb{R}^2 \mid a<x<b \wedge y=c\} = [a,b]\times\{c\} =
I_1\times\{c\}, \\[2pt]
L_2 &= \{(b,y)\in\mathbb{R}^2 \mid c<y<d\}
= \{(x,y)\in\mathbb{R}^2 \mid x=b \wedge c<y<d\} = \{b\}\times[c,d] =
\{b\}\times I_2, \\[2pt]
L_3 &= \{(x,d)\in\mathbb{R}^2 \mid a<x<b\}
= \{(x,y)\in\mathbb{R}^2 \mid a<x<b \wedge y=d\} = [a,b]\times\{d\} =
I_1\times\{d\}, \\[2pt]
L_4 &= \{(a,y)\in\mathbb{R}^2 \mid c<y<d\}
= \{(x,y)\in\mathbb{R}^2 \mid x=a \wedge c<y<d\} = \{a\}\times[c,d] =
\{a\}\times I_2 .
\end{aligned}
\]

Por lo tanto,
\[
\begin{aligned}
\partial S &= I_1\times\{c\} \;\cup\; \{b\}\times I_2 \;\cup\; I_1\times\{d\}
\;\cup\; \{a\}\times I_2 \\
&= I_1\times(\{c\}\cup\{d\}) \;\cup\; (\{a\}\cup\{b\})\times I_2 \\
&= (I_1\times \partial I_2) \;\cup\; (\partial I_1 \times I_2) \\
&= \partial(I_1\times I_2) .
\end{aligned}
\]
\end{ejemplo}

\begin{observacion}
El resultado del Ejemplo~\ref{ej:frontera-rectangulo},
\[
\partial(I_1\times I_2) = (\partial I_1\times I_2) \cup (I_1\times
\partial I_2),
\]
es particularmente importante: será usado más adelante para el formalismo
lagrangiano de objetos extendidos, donde la frontera de un dominio producto
determina las condiciones de borde del sistema físico.
\end{observacion}

\chapter{Transformación de Coordenadas}

\section{Motivación}

Al estudiar un sistema mecánico, la elección del sistema de coordenadas no
es única: un mismo punto del espacio se puede describir en coordenadas
cartesianas, cilíndricas, esféricas, o en cualquier otro sistema
conveniente para el problema. Un problema con simetría esférica (como una
fuerza central) resulta mucho más simple en coordenadas esféricas que en
cartesianas. El objetivo de este capítulo es formalizar qué significa
``cambiar de coordenadas'' y cómo transforman, bajo dicho cambio, los
distintos objetos geométricos asociados a un vector: sus componentes, la
base coordenada, la base dual y la métrica.

\section{Coordinatización del espacio}

\begin{definicion}[Coordinatización]
Sea $M$ el espacio físico (por ejemplo $\mathbb{R}^3$). Una
\textbf{coordinatización} de $M$ es un mapeo biyectivo
\[
x : M \longrightarrow \mathbb{R}^3, \qquad p \longmapsto x(p) =
(x^1,x^2,x^3),
\]
que asigna a cada punto $p \in M$ una terna ordenada de números reales,
llamados las \textbf{coordenadas} de $p$.
\end{definicion}

\begin{ejemplo}
Las coordenadas cartesianas, cilíndricas y esféricas son tres
coordinatizaciones distintas de un mismo punto $p \in \mathbb{R}^3$:
\[
x_{\text{cart}}(p) = (x,y,z), \qquad
x_{\text{cil}}(p) = (\rho,\varphi,z), \qquad
x_{\text{esf}}(p) = (r,\varphi,\theta).
\]
\end{ejemplo}

\section{Transformación de coordenadas}

\begin{definicion}[Transformación de coordenadas]
Sean $x^i$ y $x'^i$ ($i=1,2,3$) dos coordinatizaciones distintas de un
mismo punto $p$. Una \textbf{transformación de coordenadas} es un par de
mapeos biyectivos e inversos entre sí,
\[
x'^i = x'^i(x^j), \qquad x^i = x^i(x'^j),
\]
que relacionan ambas coordinatizaciones.
\end{definicion}

\begin{ejemplo}[Cartesianas a esféricas]
La transformación de coordenadas cartesianas $(x,y,z)$ a esféricas
$(r,\varphi,\theta)$ está dada por
\[
x = r\sin\theta\cos\varphi, \qquad y = r\sin\theta\sin\varphi, \qquad
z = r\cos\theta,
\]
con transformación inversa
\[
r = \sqrt{x^2+y^2+z^2}, \qquad
\varphi = \arctan\!\left(\frac{y}{x}\right), \qquad
\theta = \arccos\!\left(\frac{z}{\sqrt{x^2+y^2+z^2}}\right).
\]
\end{ejemplo}

\subsection{El jacobiano de la transformación}

\begin{definicion}[Matriz jacobiana]
Dada una transformación de coordenadas $x'^i=x'^i(x^j)$, se define la
\textbf{matriz jacobiana} de la transformación como la matriz de
$3\times 3$
\[
J = \left[ \frac{\partial x'^i}{\partial x^j} \right]_{i,j=1,2,3} =
\begin{pmatrix}
\dfrac{\partial x'^1}{\partial x^1} & \dfrac{\partial x'^1}{\partial x^2} &
\dfrac{\partial x'^1}{\partial x^3} \\[8pt]
\dfrac{\partial x'^2}{\partial x^1} & \dfrac{\partial x'^2}{\partial x^2} &
\dfrac{\partial x'^2}{\partial x^3} \\[8pt]
\dfrac{\partial x'^3}{\partial x^1} & \dfrac{\partial x'^3}{\partial x^2} &
\dfrac{\partial x'^3}{\partial x^3}
\end{pmatrix},
\]
y el \textbf{jacobiano} de la transformación es su determinante,
$J = \det\!\left[\partial x'^i/\partial x^j\right]$.
\end{definicion}

\begin{observacion}
Para que la transformación sea invertible en un entorno de un punto se
requiere $J\neq 0$ en dicho punto (teorema de la función inversa). Los puntos donde $J=0$ son puntos singulares de la coordinatización; por
ejemplo, el origen $r=0$ en coordenadas esféricas, o el eje $z$
($\rho=0$) en coordenadas cilíndricas.
\end{observacion}

\begin{ejemplo}
Para la transformación de cartesianas a esféricas, el jacobiano es
\[
J = \det\!\left[\frac{\partial x'^i}{\partial x^j}\right] = r^2\sin\theta .
\]
El jacobiano se anula en $r=0$ y en $\theta=0,\pi$ (el eje $z$), que son
precisamente los puntos donde la coordinatización esférica es singular:
en el origen $\varphi$ y $\theta$ no están definidos, y sobre el eje $z$,
$\varphi$ no está definido.
\end{ejemplo}

\section{Vectores base}

\begin{definicion}[Base coordenada]
Dada una coordinatización $x^i$ de $\mathbb{R}^3$, la \textbf{base
coordenada} asociada es el conjunto de vectores
\[
\vec{e}_i = \frac{\partial \vec{r}}{\partial x^i}, \qquad i=1,2,3,
\]
donde $\vec{r}=\vec r(x^1,x^2,x^3)$ es el vector posición.
\end{definicion}

\begin{proposicionc}
Bajo una transformación de coordenadas $x^i = x^i(x'^j)$, la base
coordenada $\{\vec e_i\}_{i=1}^3$ transforma como
\[
\vec{e}\,'_i = \sum_{j=1}^{3} \frac{\partial x^j}{\partial x'^i}\, \vec e_j
.
\]
\end{proposicionc}

\begin{proof}
Por regla de la cadena,
\[
\vec e\,'_i = \frac{\partial \vec r}{\partial x'^i} =
\sum_{j=1}^{3} \frac{\partial x^j}{\partial x'^i}
\frac{\partial \vec r}{\partial x^j} =
\sum_{j=1}^{3} \frac{\partial x^j}{\partial x'^i}\, \vec e_j . \qedhere
\]
\end{proof}

De manera análoga, contrayendo la expresión anterior con
$\sum_i \partial x'^k/\partial x^i$ y usando
$\sum_i (\partial x'^k/\partial x^i)(\partial x^i/\partial x'^j) =
\delta^k_j$, se obtiene la relación inversa:

\begin{proposicionc}
Bajo la transformación inversa $x'^i=x'^i(x^j)$,
\[
\vec e_i = \sum_{j=1}^{3} \frac{\partial x'^j}{\partial x^i}\, \vec e\,'_j
.
\]
\end{proposicionc}

\begin{ejemplo}[Base esférica a partir de la cartesiana]
Sea $\{\vec e_1,\vec e_2,\vec e_3\} = \{\ihat,\jhat,\khat\}$ la base
cartesiana, y sea $\{\vec e\,'_1,\vec e\,'_2,\vec e\,'_3\} =
\{\vec e_r,\vec e_\varphi,\vec e_\theta\}$ la base coordenada esférica.
Usando $x^j=x^j(x'^i)=x^j(r,\varphi,\theta)$,
\[
\vec e\,'_1 = \vec e_r = \sum_{j=1}^3 \frac{\partial x^j}{\partial r}
\vec e_j = \frac{\partial x}{\partial r}\ihat + \frac{\partial y}{\partial
r}\jhat + \frac{\partial z}{\partial r}\khat ,
\]
lo que da
\[
\vec e_r = \sin\theta\cos\varphi\,\ihat + \sin\theta\sin\varphi\,\jhat +
\cos\theta\,\khat .
\]
Análogamente, derivando respecto de $\varphi$ y de $\theta$,
\[
\vec e_\varphi = r\sin\theta\,(-\sin\varphi\,\ihat + \cos\varphi\,\jhat),
\qquad
\vec e_\theta = r\big(\cos\varphi\cos\theta\,\ihat +
\sin\varphi\cos\theta\,\jhat - \sin\theta\,\khat\big).
\]

\medskip
Notar que $\vec e_r$ es unitario, pero $\vec e_\varphi$ y $\vec e_\theta$
no lo son: $\|\vec e_\varphi\| = r\sin\theta$ y $\|\vec e_\theta\|=r$. La
base coordenada $\{\vec e_r,\vec e_\varphi,\vec e_\theta\}$ es ortogonal
pero \textbf{no ortonormal}; ésta es la base ``natural'' asociada a las
coordenadas esféricas, distinta de la base ortonormal
$\{\hat r,\hat\varphi,\hat\theta\}$ usada habitualmente en física
elemental.

\medskip
Realizando el cálculo inverso, $x^i=x^i(x'^j)$, se obtienen los vectores
cartesianos en términos de la base esférica:
\[
\ihat = \frac{x}{\sqrt{x^2+y^2+z^2}}\,\vec e_r - \frac{y}{x^2+y^2}\,
\vec e_\varphi + \frac{xz}{\sqrt{x^2+y^2}\,(x^2+y^2+z^2)}\, \vec e_\theta ,
\]
y expresiones análogas para $\jhat$ y $\khat$; reescritas en coordenadas
esféricas por estética,
\[
\ihat = \cos\varphi\sin\theta\,\vec e_r - \frac{\sin\varphi}{r\sin\theta}\,
\vec e_\varphi + \frac{\cos\varphi\cos\theta}{r}\,\vec e_\theta, \quad
\jhat = \sin\varphi\sin\theta\,\vec e_r + \frac{\cos\varphi}{r\sin\theta}\,
\vec e_\varphi + \frac{\sin\varphi\cos\theta}{r}\,\vec e_\theta,
\]
\[
\khat = \cos\theta\,\vec e_r - \frac{\sin\theta}{r}\,\vec e_\theta .
\]
\end{ejemplo}

\section{Componentes de un vector: tensores contravariantes}

\begin{proposicionc}
\label{prop:transf-componentes}
Sea $\vec A$ un vector con componentes $A^i$ en la base $\{\vec e_i\}$.
Bajo la transformación de coordenadas $x^i=x^i(x'^j)$, las componentes de
$\vec A$ transforman como
\[
A'^i = \sum_{j=1}^{3} \frac{\partial x'^i}{\partial x^j}\, A^j .
\]
\end{proposicionc}

\begin{proof}
Como $\vec A$ no depende del sistema coordenado, $\vec A\big|_{O'} =
\vec A\big|_O$, es decir
\[
\sum_{i=1}^3 A'^i \vec e\,'_i = \sum_{j=1}^3 A^j \vec e_j =
\sum_{j=1}^3 A^j \sum_{i=1}^3 \frac{\partial x'^i}{\partial x^j}\vec e\,'_i
= \sum_{i=1}^3 \left(\sum_{j=1}^3 \frac{\partial x'^i}{\partial x^j}
A^j\right) \vec e\,'_i .
\]
Como la base $\{\vec e\,'_i\}$ es linealmente independiente, se concluye
\[
A'^i = \sum_{j=1}^3 \frac{\partial x'^i}{\partial x^j} A^j . \qedhere
\]
\end{proof}

Contrayendo con $\sum_i \partial x^k/\partial x'^i$ y usando la regla de
la cadena se obtiene la relación inversa:

\begin{proposicionc}
\[
A^k = \sum_{i=1}^{3} \frac{\partial x^k}{\partial x'^i}\, A'^i .
\]
\end{proposicionc}

\begin{ejemplo}
Sea $\vec A = A_x\ihat + A_y\jhat + A_z\khat = A_r\vec e_r +
A_\varphi\vec e_\varphi + A_\theta\vec e_\theta$. Usando la
Proposición~\ref{prop:transf-componentes} con la transformación a
esféricas se obtiene, tras simplificar,
\[
A_r = \cos\varphi\sin\theta\,A_x + \sin\varphi\sin\theta\,A_y +
\cos\theta\,A_z = \vec e_r\cdot\vec A ,
\]
\[
A_\varphi = -\frac{\sin\varphi}{r\sin\theta}A_x +
\frac{\cos\varphi}{r\sin\theta}A_y = \frac{1}{r^2\sin^2\theta}\,
\vec e_\varphi\cdot\vec A , \qquad
A_\theta = \frac{\cos\varphi\cos\theta}{r}A_x +
\frac{\sin\varphi\cos\theta}{r}A_y - \frac{\sin\theta}{r}A_z =
\frac{1}{r^2}\,\vec e_\theta\cdot\vec A .
\]
Notar que $A_\varphi$ y $A_\theta$ \textbf{no} son simplemente
$\vec e_\varphi\cdot\vec A$ ni $\vec e_\theta\cdot\vec A$, sino que llevan
un factor de normalización: esto ocurre porque $\{\vec e_r,\vec
e_\varphi,\vec e_\theta\}$ no es una base ortonormal (véase la sección
sobre la base dual, más adelante).

La transformación inversa da
\[
A_x = \frac{x}{\sqrt{x^2+y^2+z^2}}A_r - y A_\varphi +
\frac{xz}{\sqrt{x^2+y^2}}A_\theta, \qquad
A_y = \frac{y}{\sqrt{x^2+y^2+z^2}}A_r + x A_\varphi +
\frac{yz}{\sqrt{x^2+y^2}}A_\theta,
\]
\[
A_z = \frac{z}{\sqrt{x^2+y^2+z^2}}A_r - \sqrt{x^2+y^2}\,A_\theta .
\]
\end{ejemplo}

\begin{observacion}
De la discusión anterior se concluye que:
\begin{itemize}
\item Las cantidades con el índice \textbf{arriba} transforman con la
matriz $\partial x'^i/\partial x^j$.
\item Las cantidades con el índice \textbf{abajo} transforman con la
matriz $\partial x^i/\partial x'^j$.
\end{itemize}
En la construcción de un vector $\vec A = \sum_i A^i \vec e_i$ se
necesitan las cantidades $\{\vec e_i\}$ (índice abajo) y $\{A^i\}$
(índice arriba): se contrae un índice arriba con uno abajo. Esta es la
razón estructural por la que los vectores no dependen del sistema
coordenado: las dos leyes de transformación (una y su inversa) se
cancelan exactamente.
\end{observacion}

\begin{definicion}[Índices contraídos]
Un índice arriba y un índice abajo se dicen \textbf{contraídos} si están
bajo una suma sobre dicho índice.
\end{definicion}

\begin{definicion}[Tensor contravariante de orden uno]
La cantidad $T^i$ se dice un \textbf{tensor contravariante de orden uno},
o \textbf{tensor de orden $(1,0)$}, si bajo una transformación de
coordenadas $x^i=x^i(x'^j)$ transforma como
\[
T'^i = \sum_{j=1}^{3} \frac{\partial x'^i}{\partial x^j}\, T^j .
\]
\end{definicion}

Las componentes $A^i$ de un vector arbitrario $\vec A$ son, por lo tanto,
un tensor de orden $(1,0)$.

\section{Vectores duales y tensores covariantes}

\begin{definicion}[Base dual]
La \textbf{base dual} $\{\vec E^i\}_{i=1}^3$ asociada a la base
coordenada $\{\vec e_i\}_{i=1}^3$ es el conjunto de (co)vectores definido
por la relación invariante
\[
E^i(\vec e_j) = \delta^i_j ,
\]
que se cumple en cualquier sistema de coordenadas:
$E^i(e_j)\big|_O = E^i(e_j)\big|_{O'} = \delta^i_j$.
\end{definicion}

\begin{observacion}
Los vectores duales pertenecen al espacio dual $V^*(n,K)$; un vector dual
arbitrario se escribe $\vec f = \sum_i f_i \vec E^i$, con componentes con
el índice abajo. Como $\vec f$ tampoco depende del sistema coordenado
($\vec f\big|_{O'}=\vec f\big|_O$), la base dual debe poseer su propia
ley de transformación, distinta de la de la base coordenada, ya que
$\vec E^i$ lleva el índice arriba.
\end{observacion}

\begin{proposicionc}
\label{prop:transf-base-dual}
Bajo la transformación de coordenadas $x^i=x^i(x'^j)$, la base dual
transforma como
\[
E'^i = \sum_{j=1}^{3} \frac{\partial x'^i}{\partial x^j}\, E^j .
\]
\end{proposicionc}

\begin{proof}
Partimos de la relación invariante $E^i(e_j)\big|_{O'} = \delta^i_j$ y
usamos la ley de transformación de $\vec e_j$ (transformación inversa
$x^j=x^j(x'^k)$):
\[
E'^i(e'_j) = \left(\sum_{k=1}^3 \frac{\partial x'^i}{\partial x^k} E^k
\right)\!\left(\sum_{l=1}^3 \frac{\partial x^l}{\partial x'^j} e_l\right)
= \sum_{k=1}^3\sum_{l=1}^3 \frac{\partial x'^i}{\partial x^k}
\frac{\partial x^l}{\partial x'^j}\, E^k(e_l)
= \sum_{k=1}^3 \frac{\partial x'^i}{\partial x^k}\frac{\partial
x^k}{\partial x'^j} = \delta^i_j , \qedhere
\]
lo cual confirma que la ley propuesta es consistente con la definición
invariante de la base dual.
\end{proof}

Contrayendo con $\sum_i \partial x^k/\partial x'^i$ se obtiene la
transformación inversa:
\[
E^k = \sum_{i=1}^{3} \frac{\partial x^k}{\partial x'^i}\, E'^i .
\]

\begin{proposicionc}
\label{prop:transf-componentes-dual}
La componente $f_i$ de un vector dual $\vec f$ transforma, bajo
$x'^i=x'^i(x^j)$, como
\[
f'_i = \sum_{j=1}^{3} \frac{\partial x^j}{\partial x'^i}\, f_j .
\]
\end{proposicionc}

\begin{proof}
Como $\vec f\big|_O = \vec f\big|_{O'}$,
\[
\sum_{j=1}^3 f_j E^j = \sum_{i=1}^3 f'_i E'^i = \sum_{i=1}^3 f'_i
\sum_{j=1}^3 \frac{\partial x^j}{\partial x'^i} E^j ,
\]
y por independencia lineal de $\{\vec E^j\}$,
\[
f'_i = \sum_{j=1}^3 \frac{\partial x^j}{\partial x'^i} f_j . \qedhere
\]
\end{proof}

\begin{definicion}[Tensor covariante de orden uno]
La cantidad $T_i$ se dice un \textbf{tensor covariante de orden uno}, o
\textbf{tensor de orden $(0,1)$}, si bajo la transformación de
coordenadas $x^i=x^i(x'^j)$ transforma como
\[
T'_i = \sum_{j=1}^{3} \frac{\partial x^j}{\partial x'^i}\, T_j .
\]
\end{definicion}

\begin{observacion}
El término \textbf{covarianza} (invariancia de forma) de un vector
significa que el vector no cambia de forma bajo una transformación de
coordenadas: no cambia su dirección, sentido ni magnitud. Sin embargo,
bajo un cambio de coordenadas cambian las bases coordenadas, por lo que
el valor numérico de sus componentes sí cambia.
\end{observacion}

\section{La métrica}

\begin{definicion}[Métrica]
La \textbf{métrica} asociada a una base coordenada $\{\vec e_i\}_{i=1}^3$
es
\[
g_{ij} = \vec e_i \cdot \vec e_j .
\]
\end{definicion}

\begin{ejemplo}
En coordenadas cartesianas, $g_{ij}=\delta_{ij}$ (la \textbf{métrica de
Euclides}). En coordenadas esféricas,
\[
g_{ij} = \begin{pmatrix} 1 & 0 & 0 \\ 0 & r^2\sin^2\theta & 0 \\ 0 & 0 &
r^2 \end{pmatrix},
\]
consistente con $\|\vec e_r\|^2=1$, $\|\vec e_\varphi\|^2 = r^2\sin^2
\theta$, $\|\vec e_\theta\|^2=r^2$ obtenidas anteriormente.
\end{ejemplo}

\begin{proposicionc}
\label{prop:transf-metrica}
La métrica $g_{ij}(x^k)$ bajo una transformación de coordenadas
$x^i=x^i(x'^j)$ transforma como
\[
g'_{ij} = \sum_{k=1}^{3}\sum_{l=1}^{3} \frac{\partial x^k}{\partial x'^i}
\frac{\partial x^l}{\partial x'^j}\, g_{kl} .
\]
\end{proposicionc}

\begin{proof}
Usando la ley de transformación de la base coordenada,
\[
g'_{ij} = \vec e\,'_i \cdot \vec e\,'_j = \left(\sum_{k=1}^3
\frac{\partial x^k}{\partial x'^i}\vec e_k\right)\!\cdot\!\left(
\sum_{l=1}^3 \frac{\partial x^l}{\partial x'^j}\vec e_l\right)
= \sum_{k=1}^3\sum_{l=1}^3 \frac{\partial x^k}{\partial x'^i}
\frac{\partial x^l}{\partial x'^j}\, \vec e_k\cdot\vec e_l =
\sum_{k=1}^3\sum_{l=1}^3 \frac{\partial x^k}{\partial x'^i}\frac{\partial
x^l}{\partial x'^j} g_{kl} . \qedhere
\]
\end{proof}

\begin{definicion}[Tensor covariante de orden dos]
La cantidad $T_{ij}$ se dice un tensor covariante de orden $2$, o tensor
de orden $(0,2)$, si transforma como
\[
T'_{ij} = \sum_{k=1}^{3}\sum_{l=1}^{3} \frac{\partial x^k}{\partial
x'^i}\frac{\partial x^l}{\partial x'^j}\, T_{kl} .
\]
\end{definicion}

Por lo tanto, la métrica $g_{ij}$ es un tensor de orden $(0,2)$.

\begin{ejemplo}
Recuperemos la métrica esférica a partir de la métrica cartesiana
$g_{ij}=\delta_{ij}$ usando la Proposición~\ref{prop:transf-metrica}:
\[
g'_{ij} = \sum_{k=1}^3 \frac{\partial x^k}{\partial x'^i}\frac{\partial
x^k}{\partial x'^j} = \frac{\partial x}{\partial x'^i}\frac{\partial
x}{\partial x'^j} + \frac{\partial y}{\partial x'^i}\frac{\partial
y}{\partial x'^j} + \frac{\partial z}{\partial x'^i}\frac{\partial
z}{\partial x'^j} .
\]
Sustituyendo $x=r\sin\theta\cos\varphi$, etc., se recupera
\[
g'_{ij} = \begin{pmatrix} 1 & 0 & 0 \\ 0 & r^2\sin^2\theta & 0 \\ 0 & 0 &
r^2 \end{pmatrix}.
\]
\end{ejemplo}

\subsection{El elemento de línea}

\begin{proposicionc}
El elemento de línea $ds$, definido por $ds^2 = \sum_{i,j} g_{ij}\,
dx^i\,dx^j$, es un invariante bajo transformaciones de coordenadas:
$ds' = ds$.
\end{proposicionc}

\begin{proof}
En el sistema $O'x'^i$,
\[
ds'^2 = \sum_{i=1}^3\sum_{j=1}^3 g'_{ij}\, dx'^i\, dx'^j .
\]
Bajo $x'^i=x'^i(x^j)$, el diferencial $dx'^i$ transforma como
$dx'^i = \sum_j (\partial x'^i/\partial x^j)\, dx^j$, es decir, $dx^i$ es
un tensor de orden $(1,0)$. Sustituyendo esto y la ley de transformación
de $g'_{ij}$ (Proposición~\ref{prop:transf-metrica}),
\[
ds'^2 = \sum_{i,j,k,l} \frac{\partial x^k}{\partial x'^i}\frac{\partial
x^l}{\partial x'^j}\, g_{kl} \sum_{n} \frac{\partial x'^i}{\partial x^n}
dx^n \sum_{m}\frac{\partial x'^j}{\partial x^m} dx^m
= \sum_{k,l,n,m} \left(\sum_i \frac{\partial x^k}{\partial x'^i}
\frac{\partial x'^i}{\partial x^n}\right)\!\!\left(\sum_j \frac{\partial
x^l}{\partial x'^j}\frac{\partial x'^j}{\partial x^m}\right) g_{kl}\,
dx^n dx^m .
\]
Ambos paréntesis son deltas de Kronecker, $\delta^k_n$ y $\delta^l_m$, de
modo que
\[
ds'^2 = \sum_{k,l} g_{kl}\, dx^k dx^l = ds^2 \quad\Longrightarrow\quad
ds' = ds . \qedhere
\]
\end{proof}

\begin{observacion}
El elemento de línea no es un vector sino un \emph{escalar}: un número
(con dimensión) cuyo valor no cambia bajo un cambio de coordenadas. En
particular, el largo de una trayectoria tiene el mismo valor en todos los
sistemas coordenados.
\end{observacion}

\section{Invariantes: producto escalar, magnitud y ángulo}

\begin{proposicionc}
El producto escalar $\vec A\cdot\vec B = \sum_{i,j} g_{ij} A^i B^j$ entre
dos vectores cualesquiera es invariante bajo transformación de
coordenadas.
\end{proposicionc}

\begin{proof}
Análogo a la demostración de la invariancia del elemento de línea,
reemplazando $dx^i$ por $A^i$ y $dx^j$ por $B^j$ (ambos tensores de orden
$(1,0)$):
\[
\vec A\cdot\vec B\big|_{O'} = \sum_{i,j} g'_{ij} A'^i B'^j = \sum_{n,m}
g_{nm} A^n B^m = \vec A\cdot\vec B\big|_O . \qedhere
\]
\end{proof}

El producto escalar es un número real, no un vector; por lo tanto su
valor numérico no cambia bajo una transformación de coordenadas. De este
resultado se derivan dos corolarios inmediatos.

\begin{corolario}
La magnitud $\|\vec A\| = \sqrt{\vec A\cdot\vec A}$ de un vector es
invariante bajo transformaciones de coordenadas.
\end{corolario}

\begin{proof}
Tomando $\vec B=\vec A$ en la proposición anterior,
$\vec A\cdot\vec A\big|_{O'} = \vec A\cdot\vec A\big|_O$, es decir
$\|\vec A\|^2\big|_{O'} = \|\vec A\|^2\big|_O$, luego
$\|\vec A\|\big|_{O'} = \|\vec A\|\big|_O$. \qedhere
\end{proof}

\begin{corolario}
El ángulo entre dos vectores es invariante bajo transformación de
coordenadas.
\end{corolario}

\begin{proof}
El ángulo $\theta$ entre $\vec A$ y $\vec B$ está dado por
$\cos\theta = (\vec A\cdot\vec B)/(\|\vec A\|\|\vec B\|)$. Como el
numerador y ambos factores del denominador son invariantes,
$\cos\theta\big|_{O'} = \cos\theta\big|_O$, y por lo tanto
$\theta\big|_{O'} = \theta\big|_O$. \qedhere
\end{proof}

\begin{observacion}[Definición formal vía el espacio dual]
Formalmente, el producto escalar entre un vector dual $\vec f = \sum_i
f_i \vec E^i \in V^*(n,K)$ y un vector $\vec v = \sum_i v^i \vec e_i \in
V(n,K)$ se define como el emparejamiento natural
\[
f(v) = \sum_{i=1}^{3} f_i\, v^i .
\]
\end{observacion}

\begin{proposicionc}
El emparejamiento $f(v) = \sum_i f_i v^i$ es invariante bajo la
transformación de coordenadas $x'^i=x'^i(x^j)$.
\end{proposicionc}

\begin{proof}
Las cantidades $v^i$ y $f_i$ son tensores de orden $(1,0)$ y $(0,1)$
respectivamente, es decir $v'^i = \sum_j (\partial x'^i/\partial x^j)
v^j$ y $f'_i = \sum_j (\partial x^j/\partial x'^i) f_j$. Entonces
\[
f'(v') = \sum_{i=1}^3 f'_i v'^i = \sum_{i,j,k}\frac{\partial
x^j}{\partial x'^i} f_j\, \frac{\partial x'^i}{\partial x^k} v^k =
\sum_{j,k}\left(\sum_i \frac{\partial x^j}{\partial x'^i}\frac{\partial
x'^i}{\partial x^k}\right) f_j v^k = \sum_{j,k} \delta^j_k\, f_j v^k =
\sum_{j} f_j v^j = f(v). \qedhere
\]
\end{proof}

\newpage
\section{La métrica inversa}

\begin{proposicionc}
La métrica inversa $g^{ij}$ está dada en términos de la base dual
$\{\vec E^i\}_{i=1}^3$ por
\[
g^{ij} = \vec E^i \cdot \vec E^j .
\]
\end{proposicionc}

\begin{proof}
La base dual se relaciona con la base coordenada a través de la métrica
inversa, $\vec E^i = \sum_j g^{ij} \vec e_j$, de modo que se cumpla la
definición $E^i(e_j) = \vec E^i\cdot\vec e_j = \sum_k g^{ik}\vec
e_k\cdot\vec e_j = \sum_k g^{ik} g_{kj} = \delta^i_j$. Entonces
\[
\vec E^i\cdot\vec E^j = \sum_{k} g^{ik}\vec e_k \cdot \sum_{l} g^{jl}\vec
e_l = \sum_{k,l} g^{ik} g^{jl}\, g_{kl} = \sum_k g^{ik}\left(\sum_l
g^{jl}g_{kl}\right) = \sum_k g^{ik}\, \delta^j_k = g^{ij} . \qedhere
\]
\end{proof}

\begin{ejemplo}
En coordenadas esféricas, usando $\vec E^1=\vec e_r$, $\vec E^2 = \vec
e_\varphi/(r^2\sin^2\theta)$, $\vec E^3 = \vec e_\theta/r^2$ (obtenidos de
$\vec E^i = \sum_j g^{ij}\vec e_j$ con $g^{ij}=\text{diag}(1,\,
1/(r^2\sin^2\theta),\,1/r^2)$), se verifica directamente
\[
g^{11} = \vec E^1\cdot\vec E^1 = \vec e_r\cdot\vec e_r = 1, \qquad
g^{22} = \vec E^2\cdot\vec E^2 = \frac{\vec e_\varphi\cdot\vec
e_\varphi}{r^4\sin^4\theta} = \frac{1}{r^2\sin^2\theta}, \qquad
g^{33} = \vec E^3\cdot\vec E^3 = \frac{\vec e_\theta\cdot\vec
e_\theta}{r^4} = \frac{1}{r^2},
\]
y los términos cruzados $g^{12}=g^{13}=g^{23}=0$ por ortogonalidad de la
base esférica. Es decir,
\[
g^{ij} = \begin{pmatrix} 1 & 0 & 0 \\ 0 & \dfrac{1}{r^2\sin^2\theta} & 0
\\ 0 & 0 & \dfrac{1}{r^2} \end{pmatrix},
\]
el resultado conocido.
\end{ejemplo}

\begin{proposicionc}
\label{prop:transf-metrica-inversa}
La métrica inversa $g^{ij}$ bajo una transformación de coordenadas
$x'^i=x'^i(x^j)$ transforma como
\[
g'^{ij} = \sum_{k=1}^{3}\sum_{l=1}^{3} \frac{\partial x'^i}{\partial
x^k}\frac{\partial x'^j}{\partial x^l}\, g^{kl} .
\]
\end{proposicionc}

\begin{proof}
Escribiendo $g'^{ij} = \vec E'^i\cdot\vec E'^j$ y usando la ley de
transformación de la base dual (Proposición~\ref{prop:transf-base-dual}),
\[
g'^{ij} = \left(\sum_k \frac{\partial x'^i}{\partial x^k}\vec
E^k\right)\!\cdot\!\left(\sum_l \frac{\partial x'^j}{\partial x^l}\vec
E^l\right) = \sum_{k,l} \frac{\partial x'^i}{\partial x^k}\frac{\partial
x'^j}{\partial x^l}\, \vec E^k\cdot\vec E^l = \sum_{k,l} \frac{\partial
x'^i}{\partial x^k}\frac{\partial x'^j}{\partial x^l}\, g^{kl} .
\qedhere
\]
\end{proof}

\begin{ejemplo}
Verifiquemos la componente $g'^{11}$ de la métrica esférica a partir de
$g^{kl}=\delta^{kl}$ (cartesiana) usando la
Proposición~\ref{prop:transf-metrica-inversa}:
\[
g'^{11} = \left(\frac{\partial r}{\partial x}\right)^2 +
\left(\frac{\partial r}{\partial y}\right)^2 + \left(\frac{\partial
r}{\partial z}\right)^2 = \frac{x^2}{r^2} + \frac{y^2}{r^2} +
\frac{z^2}{r^2} = 1 ,
\]
consistente con el resultado directo. Repitiendo el cálculo para las
demás componentes se recupera $g^{ij} =
\text{diag}(1,\,1/(r^2\sin^2\theta),\,1/r^2)$.
\end{ejemplo}

\begin{definicion}[Tensor contravariante de orden dos]
La cantidad $T^{ij}$ se dice un tensor contravariante de orden $2$, o
tensor de orden $(2,0)$, si bajo $x'^i=x'^i(x^j)$ transforma como
\[
T'^{ij} = \sum_{k=1}^{3}\sum_{l=1}^{3} \frac{\partial x'^i}{\partial
x^k}\frac{\partial x'^j}{\partial x^l}\, T^{kl} .
\]
\end{definicion}

Por lo tanto, la métrica inversa $g^{ij}$ es un tensor de orden $(2,0)$.

\section{Los tres símbolos delta}

Conviene distinguir con cuidado tres cantidades que, a primera vista,
parecen representar todas a la matriz identidad: $\delta_{ij}$,
$\delta^{ij}$ y $\delta^i_j$.

\begin{proposicionc}
De las tres cantidades $\delta_{ij}$, $\delta^{ij}$ y $\delta^i_j$, solo
$\delta^i_j$ es invariante bajo transformaciones de coordenadas.
\end{proposicionc}

\begin{proof}
Como $\delta^k_l$ es un tensor de orden $(1,1)$ (transforma con un
índice arriba y uno abajo, cada uno con su matriz correspondiente),
\[
\delta'^i_j = \sum_{k=1}^3\sum_{l=1}^3 \frac{\partial x'^i}{\partial
x^k}\frac{\partial x^l}{\partial x'^j}\, \delta^k_l = \sum_{k=1}^3
\frac{\partial x'^i}{\partial x^k}\frac{\partial x^k}{\partial x'^j} =
\frac{\partial x'^i}{\partial x'^j} = \delta^i_j . \qedhere
\]
\end{proof}

En cambio, $\delta_{ij}$ y $\delta^{ij}$ transforman como tensores de
orden $(0,2)$ y $(2,0)$ respectivamente, y por lo tanto \textbf{no} son
invariantes: cambian su forma funcional bajo un cambio de coordenadas
(por ejemplo, al pasar a esféricas dejan de ser la matriz identidad).

\begin{observacion}
En resumen:
\begin{itemize}
\item $\delta^i_j$ es la matriz identidad (invariante, tensor de orden
$(1,1)$).
\item $\delta_{ij}$ es la métrica asociada a las coordenadas cartesianas,
llamada \textbf{métrica de Euclides} (tensor de orden $(0,2)$).
\item $\delta^{ij}$ es la métrica inversa de la métrica de Euclides
(tensor de orden $(2,0)$).
\end{itemize}
\end{observacion}

\chapter{Análisis Tensorial}

\section{Convención de suma de Einstein}

Sea el espacio vectorial $V(n,K)$, generalización a $n$ dimensiones sobre
un campo $K$ del espacio vectorial estudiado en el capítulo anterior, con
base $\{\vec e_i\}_{i=1}^n$. Si $\vec v\in V(n,K)$, se escribe como
combinación lineal de la base elegida,
\[
\vec v = \sum_{i=1}^{n} v^i \vec e_i,
\]
donde $v^i$ es la componente del vector respecto de la base.

\begin{definicion}[Convención de Einstein]
Debido a que al trabajar con muchos índices aparecen muchas sumas,
adoptamos la \textbf{convención de suma de Einstein}: eliminaremos el
símbolo de suma $\sum$ y asumiremos que \textbf{índices repetidos
significan suma}, a menos que se diga explícitamente lo contrario.
\end{definicion}

Con esta convención, la ecuación anterior se simplifica a
\[
\vec v = v^i \vec e_i .
\]
Del mismo modo, si $\vec f \in V^*(n,K)$ es un vector dual escrito en la
base dual $\{\vec E^i\}_{i=1}^n$,
\[
\vec f = f_i \vec E^i .
\]
Los vectores del espacio dual son mapeos lineales que actúan sobre los
vectores del espacio vectorial y los llevan a $K$,
\[
f(v) = f_i v^i \in K,
\]
donde recordamos que los vectores duales se definen mediante
$E^i(e_j)=\delta^i_j$. A partir de aquí, salvo que se indique lo
contrario, se usará siempre la convención de Einstein.

\section{Subir y bajar índices}

Como se estudió en el capítulo anterior, la métrica $g_{ij}$ permite
definir el producto escalar y, además, permite pasar del espacio
vectorial al espacio vectorial dual,
\[
g_{ij} : V(n,K) \longrightarrow V^*(n,K), \qquad v_i = g_{ij} v^j .
\]
Como $\det(g_{ij})\neq 0$, existe la métrica inversa $g^{ij}$, tal que
$g^{ik}g_{kj}=\delta^i_j$, la cual permite pasar del espacio vectorial
dual al espacio vectorial,
\[
g^{ij} : V^*(n,K) \longrightarrow V(n,K), \qquad f^i = g^{ij} f_j .
\]

\begin{observacion}
Decimos que la métrica \textbf{baja índices}: convierte índices
contravariantes en covariantes. La métrica inversa \textbf{sube
índices}: convierte índices covariantes en contravariantes.
\end{observacion}

\begin{proposicionc}
La cantidad $v_i = g_{ij}v^j$, construida a partir del tensor $v^i$ de
orden $(1,0)$, es un tensor de orden $(0,1)$.
\end{proposicionc}

\begin{proof}
Sea la transformación de coordenadas $x'^i=x'^i(x^j)$. Usando la
definición $v_i=g_{ij}v^j$ en el sistema $O'x'^i$ y las leyes de
transformación de $g_{ij}$ y $v^j$,
\[
v'_i = g'_{ij} v'^j = \frac{\partial x^k}{\partial x'^i}\frac{\partial
x^l}{\partial x'^j} g_{kl}\, \frac{\partial x'^j}{\partial x^n} v^n =
\frac{\partial x^k}{\partial x'^i}\left(\frac{\partial x^l}{\partial
x'^j}\frac{\partial x'^j}{\partial x^n}\right) g_{kl} v^n = \frac{\partial
x^k}{\partial x'^i}\, \delta^l_n\, g_{kl} v^n = \frac{\partial
x^k}{\partial x'^i}\, g_{kl} v^l = \frac{\partial x^k}{\partial x'^i}
v_k,
\]
donde en el último paso se usó nuevamente $v_k=g_{kl}v^l$. Por lo tanto
\[
v'_i = \frac{\partial x^k}{\partial x'^i}\, v_k ,
\]
que es precisamente la ley de transformación de un tensor de orden
$(0,1)$. \qedhere
\end{proof}

\begin{proposicionc}
La cantidad $f^i=g^{ij}f_j$, construida a partir del tensor $f_i$ de
orden $(0,1)$, es un tensor de orden $(1,0)$.
\end{proposicionc}

\begin{proof}
Análogamente,
\[
f'^i = g'^{ij}f'_j = \frac{\partial x'^i}{\partial x^k}\frac{\partial
x'^j}{\partial x^l} g^{kl}\, \frac{\partial x^n}{\partial x'^j} f_n =
\frac{\partial x'^i}{\partial x^k}\left(\frac{\partial x'^j}{\partial
x^l}\frac{\partial x^n}{\partial x'^j}\right) g^{kl} f_n = \frac{\partial
x'^i}{\partial x^k}\, \delta^n_l\, g^{kl} f_n = \frac{\partial
x'^i}{\partial x^k}\, g^{kl}f_l = \frac{\partial x'^i}{\partial x^k}
f^k , \qedhere
\]
que es la ley de transformación de un tensor de orden $(1,0)$.
\end{proof}

\section{El producto tensorial}

\subsection{Motivación: producto tensorial de matrices columna}

Comencemos con $n=3$ y el espacio vectorial $\mathbb{R}^3$, escribiendo
los vectores como matrices columna.

\begin{definicion}[Producto tensorial de matrices columna]
Sean los vectores $v=(v^1,v^2,v^3)^T$ y $w=(w^1,w^2,w^3)^T$. Se define el
\textbf{producto tensorial} entre $v$ y $w$, denotado $\otimes$, como la
matriz
\[
M = v\otimes w = \begin{pmatrix} v^1 \\ v^2 \\ v^3 \end{pmatrix} \otimes
\begin{pmatrix} w^1 \\ w^2 \\ w^3 \end{pmatrix} = \begin{pmatrix}
v^1w^1 & v^1w^2 & v^1w^3 \\ v^2w^1 & v^2w^2 & v^2w^3 \\ v^3w^1 & v^3w^2 &
v^3w^3 \end{pmatrix}.
\]
\end{definicion}

\begin{ejemplo}
Sean $v=(1,-2,3)^T$ y $w=(1,1,2)^T$; entonces
\[
M = v\otimes w = \begin{pmatrix} 1 & 1 & 2 \\ -2 & -2 & -4 \\ 3 & 3 & 6
\end{pmatrix}.
\]
\end{ejemplo}

El producto tensorial entre dos matrices de $3\times 1$ genera una matriz
de $3\times 3$; como las matrices de $3\times 3$ forman un espacio
vectorial de dimensión $3\times 3=9$, su base canónica debe contener $9$
matrices linealmente independientes, dadas por
\[
e_{ij} = \left(\text{matriz con un } 1 \text{ en la fila } i,\ \text{columna } j,\ \text{y } 0
\text{ en el resto}\right), \qquad i,j=1,2,3.
\]
Cualquier matriz $M$ de $3\times 3$ se escribe entonces como combinación
lineal de estas 9 matrices base. Reescribiendo el producto tensorial
(11) en esta base:
\[
M = v\otimes w = v^1w^1e_{11} + v^1w^2e_{12} + \dots + v^3w^3e_{33} =
v^iw^j e_{ij} \qquad \text{(convención de Einstein)}.
\]
Por otro lado, $M$ se escribe en términos de la misma base como
$M = M^{ij}e_{ij}$, donde $M^{ij}$ es la componente de $M$ en la base
$\{e_{ij}\}$. Comparando ambas expresiones,
\[
M^{ij} = v^iw^j .
\]

\begin{observacion}
No todas las matrices de $3\times 3$ se pueden escribir en la forma
$M^{ij}=v^iw^j$: por ejemplo, la matriz
$\begin{pmatrix} 0 & -1 & 0 \\ 1 & 0 & 0 \\ 0 & 0 & 1\end{pmatrix}$
no admite vectores $v,w$ tales que $M=v\otimes w$. Sin embargo,
\textbf{todas} las matrices de $3\times3$, deriven o no de un producto
tensorial, se pueden escribir en la base $\{e_{ij}\}$; esta forma general
es la que se generaliza a continuación.
\end{observacion}

\subsection{Producto tensorial de espacios vectoriales genéricos}

El producto tensorial es un producto lineal:
\[
v\otimes(w+u) = v\otimes w + v\otimes u, \qquad (w+u)\otimes v = w\otimes
v + u\otimes v .
\]
Para $v,w\in V(n,K)$, el producto tensorial se escribe entonces
\[
M = v\otimes w = v^i\vec e_i \otimes w^j\vec e_j = v^iw^j\, \vec
e_i\otimes\vec e_j = M^{ij}\, \vec e_i\otimes\vec e_j, \qquad M^{ij}=
v^iw^j,
\]
y, comparando con el caso de matrices columna, se concluye que
\[
e_{ij} = \vec e_i \otimes \vec e_j, \qquad i,j=1,\dots,n.
\]

\begin{definicion}[Componente irreducible]
La componente $M^{ij}$ de la matriz $M$ se dice \textbf{irreducible} si
\emph{no} se puede escribir en la forma $M^{ij}=v^iw^j$, es decir, si
$M \neq v\otimes w$.
\end{definicion}

\begin{definicion}[Espacio $V(n,K)\otimes V(n,K)$]
Se define el espacio vectorial $V(n,K)\otimes V(n,K)$, con base
$\{\vec e_i\otimes\vec e_j\}_{i,j=1}^n$, tal que un vector $T$ de este
espacio se escribe como
\[
T = T^{ij}\, \vec e_i \otimes \vec e_j ,
\]
con $T^{ij}$ la componente de $T$ en la base elegida.
\end{definicion}

Bajo una transformación de coordenadas $x'^i=x'^i(x^j)$, la base
$\{\vec e_i\}$ transforma como $\vec e\,'_i = (\partial x^j/\partial
x'^i)\, \vec e_j$. Un vector $T\in V(n,K)\otimes V(n,K)$ escrito en los
sistemas $Ox^i$ y $O'x'^i$ debe satisfacer el \textbf{principio de
covariancia}, $T\big|_{O'}=T\big|_O$:

\begin{proposicionc}
Bajo la transformación de coordenadas $x'^i=x'^i(x^j)$, la componente de
un vector $T\in V(n,K)\otimes V(n,K)$ transforma como
\[
T'^{ij} = \frac{\partial x'^i}{\partial x^k}\frac{\partial
x'^j}{\partial x^l}\, T^{kl} .
\]
Es decir, $T^{ij}$ es un tensor de orden $(2,0)$.
\end{proposicionc}

\begin{proof}
De $T^{ij}\vec e_i\otimes\vec e_j = T'^{kl}\vec e\,'_k\otimes\vec e\,'_l
= T'^{kl}\dfrac{\partial x^i}{\partial x'^k}\dfrac{\partial
x^j}{\partial x'^l}\, \vec e_i\otimes\vec e_j$ y la independencia lineal
de la base $\{\vec e_i\otimes\vec e_j\}$, se obtiene
\[
T^{ij} = \frac{\partial x^i}{\partial x'^k}\frac{\partial
x^j}{\partial x'^l}\, T'^{kl},
\]
de donde, despejando $T'^{ij}$ (transformación inversa), se sigue el
resultado. \qedhere
\end{proof}

\subsection{Generalización a tensores de orden \texorpdfstring{$(p,0)$}
{(p,0)}}

\begin{definicion}
Se define el espacio vectorial $\underbrace{V(n,K)\otimes\dots\otimes
V(n,K)}_{p\text{ veces}}$, con base $\{\vec e_{i_1}\otimes\dots\otimes
\vec e_{i_p}\}_{i_1,\dots,i_p=1}^n$, tal que un vector $T$ de este
espacio se escribe
\[
T = T^{i_1i_2\dots i_p}\, \vec e_{i_1}\otimes\vec e_{i_2}\otimes\dots
\otimes\vec e_{i_p}.
\]
\end{definicion}

\begin{proposicionc}
Bajo $x'^i=x'^i(x^j)$, la componente $T^{i_1\dots i_p}$ transforma como
\[
T'^{i_1i_2\dots i_p} = \frac{\partial x'^{i_1}}{\partial x^{j_1}}
\frac{\partial x'^{i_2}}{\partial x^{j_2}} \cdots \frac{\partial
x'^{i_p}}{\partial x^{j_p}}\, T^{j_1j_2\dots j_p}.
\]
\end{proposicionc}

\begin{definicion}[Tensor contravariante de orden $p$]
La cantidad $T^{i_1\dots i_p}$ se dice un \textbf{tensor contravariante
de orden $p$}, o tensor de orden $(p,0)$, si bajo una transformación de
coordenadas transforma como en la proposición anterior. Su ley de
transformación contiene $p$ matrices $\partial x'^i/\partial x^j$ de
forma homogénea (todas del mismo tipo).
\end{definicion}

\subsection{Generalización a tensores de orden \texorpdfstring{$(0,q)$}
{(0,q)}}

De manera análoga, usando el espacio vectorial dual $V^*(n,K)$: para
$V,W\in V^*(n,K)$,
\[
T = V\otimes W = V_iE^i \otimes W_jE^j = V_iW_j\, E^i\otimes E^j =:
T_{ij}\, E^i\otimes E^j, \qquad T_{ij}=V_iW_j.
\]
Nótese que $T\notin V^*(n,K)$, sino que pertenece a un nuevo espacio
vectorial dual, y en general $T_{ij}\neq V_iW_j$ para un tensor
arbitrario.

\begin{definicion}
Se define el espacio vectorial $\underbrace{V^*(n,K)\otimes\dots\otimes
V^*(n,K)}_{q\text{ veces}}$, con base $\{E^{i_1}\otimes\dots\otimes
E^{i_q}\}_{i_1,\dots,i_q=1}^n$, tal que $T$ en este espacio se escribe
\[
T = T_{i_1i_2\dots i_q}\, E^{i_1}\otimes E^{i_2}\otimes\dots\otimes
E^{i_q}.
\]
\end{definicion}

\begin{proposicionc}
Bajo $x'^i=x'^i(x^j)$, la componente $T_{i_1\dots i_q}$ transforma como
\[
T'_{i_1i_2\dots i_q} = \frac{\partial x^{j_1}}{\partial x'^{i_1}}
\frac{\partial x^{j_2}}{\partial x'^{i_2}} \cdots \frac{\partial
x^{j_q}}{\partial x'^{i_q}}\, T_{j_1j_2\dots j_q}.
\]
\end{proposicionc}

\begin{definicion}[Tensor covariante de orden $q$]
La cantidad $T_{i_1\dots i_q}$ se dice un \textbf{tensor covariante de
orden $q$}, o tensor de orden $(0,q)$, si transforma como en la
proposición anterior; su ley de transformación contiene $q$ matrices
$\partial x^j/\partial x'^i$ de forma homogénea.
\end{definicion}

Puesto que $V^*(n,K)\otimes\dots\otimes V^*(n,K)$ (con $q$ factores) es
un espacio vectorial dual, actúa sobre $V(n,K)\otimes\dots\otimes
V(n,K)$ ($p$ factores, con $p=q$, sin índices libres): si
$V\in V(n,K)^{\otimes q}$ y $T\in V^*(n,K)^{\otimes q}$,
\[
T(V) = T_{i_1\dots i_q}\, V^{i_1\dots i_q} \in K .
\]

\subsection{Generalización a tensores de orden \texorpdfstring{$(1,1)$}
{(1,1)} y \texorpdfstring{$(p,q)$}{(p,q)}}

Para $V\in V(n,K)$ y $W\in V^*(n,K)$,
\[
T = V\otimes W = V^i\vec e_i\otimes W_jE^j = V^iW_j\, \vec e_i\otimes E^j
=: T^i_{\ j}\, \vec e_i\otimes E^j, \qquad T^i_{\ j}=V^iW_j,
\]
donde $T$ pertenece a un nuevo \textbf{espacio vectorial mixto}
$V(n,K)\otimes V^*(n,K)$ (en general $T^i_{\ j}\neq V^iW_j$).

\begin{proposicionc}
Si $T\in V(n,K)\otimes V^*(n,K)$, entonces bajo $x'^i=x'^i(x^j)$ la
componente $T^i_{\ j}$ transforma como
\[
T'^i_{\ j} = \frac{\partial x'^i}{\partial x^k}\frac{\partial
x^l}{\partial x'^j}\, T^k_{\ l} .
\]
\end{proposicionc}

\begin{definicion}[Tensor de orden $(1,1)$]
La cantidad $T^i_{\ j}$ se dice un tensor de orden $(1,1)$
(contravariante de orden uno y covariante de orden uno) si transforma
como en la proposición anterior.
\end{definicion}

\begin{observacion}
El espacio mixto se escribe convencionalmente $V(n,K)\otimes V^*(n,K)$ y
no al revés: como $V^*(n,K)$ es un espacio de funcionales lineales que
actúan sobre $V(n,K)$, escribir $V^*(n,K)\otimes V(n,K)$ podría sugerir
erróneamente que el dual actúa sobre el factor $V(n,K)$. (Alguna
literatura usa el orden opuesto; no debe generar confusión.) En la ley de
transformación de un tensor $(1,1)$ aparecen ambos tipos de matrices,
$\partial x'^i/\partial x^j$ y $\partial x^j/\partial x'^i$, pero
\emph{no} contraídas entre sí, pues actúan sobre índices distintos.
\end{observacion}

La generalización a un tensor de orden $(p,q)$ arbitrario es directa:

\begin{definicion}[Tensor de orden $(p,q)$]
Se define el espacio vectorial
\[
\underbrace{V(n,K)\otimes\dots\otimes V(n,K)}_{p\text{ veces}} \otimes
\underbrace{V^*(n,K)\otimes\dots\otimes V^*(n,K)}_{q\text{ veces}},
\]
con base $\{\vec e_{i_1}\otimes\dots\otimes\vec e_{i_p}\otimes
E^{j_1}\otimes\dots\otimes E^{j_q}\}$, tal que $T$ en este espacio se
escribe
\[
T = T^{i_1\dots i_p}_{\quad\ j_1\dots j_q}\, \vec e_{i_1}\otimes\dots
\otimes\vec e_{i_p}\otimes E^{j_1}\otimes\dots\otimes E^{j_q}.
\]
La cantidad $T^{i_1\dots i_p}_{\quad\ j_1\dots j_q}$ se dice un
\textbf{tensor de orden $(p,q)$} (contravariante de orden $p$ y
covariante de orden $q$) si bajo $x'^i=x'^i(x^j)$ transforma como
\[
T'^{i_1\dots i_p}_{\quad\ j_1\dots j_q} = \frac{\partial
x'^{i_1}}{\partial x^{k_1}}\cdots\frac{\partial x'^{i_p}}{\partial
x^{k_p}}\, \frac{\partial x^{l_1}}{\partial x'^{j_1}}\cdots\frac{\partial
x^{l_q}}{\partial x'^{j_q}}\, T^{k_1\dots k_p}_{\quad\ l_1\dots l_q},
\]
con $p$ matrices $\partial x'^i/\partial x^j$ y $q$ matrices $\partial
x^j/\partial x'^i$.
\end{definicion}

\begin{observacion}[Definición cualitativa]
Un tensor es una entidad algebraica de varias componentes que generaliza
los conceptos de escalar, vector y matriz de forma independiente del
sistema de coordenadas elegido: es \textbf{invariante en forma} ante un
cambio de coordenadas, $T\big|_O = T\big|_{O'}$.
\end{observacion}

\section{Subir y bajar índices en tensores generales}

La métrica, al actuar sobre un tensor de orden $(0,p)$ (o, en general,
sobre cualquier índice contravariante de un tensor de orden $(p,q)$),
genera la cantidad
\[
T^{i_1\dots i_{s-1}}_{\quad\quad\ \ i_{s+1}\dots i_p\, j} = g_{jk}\,
T^{i_1\dots i_{s-1}k i_{s+1}\dots i_p},
\]
que es un tensor de orden $(p-1,1)$ (se ``bajó'' un índice). En general
se pueden bajar $r$ índices con $r$ métricas:
\[
T^{i_1\dots i_s}_{\quad\ \ i_{s+r+1}\dots i_p\, k_1\dots k_r} =
g_{k_1j_1}\cdots g_{k_rj_r}\, T^{i_1\dots i_s j_1\dots j_r
i_{s+r+1}\dots i_p},
\]
que es un tensor de orden $(p-r,r)$. De modo análogo, la métrica inversa
permite \textbf{subir} $r$ índices covariantes de un tensor de orden
$(0,q)$, obteniendo un tensor de orden $(r,q-r)$.

\subsection{Los símbolos de Christoffel}

\begin{definicion}[Símbolos de Christoffel]
Los \textbf{símbolos de Christoffel} $\Gamma^k_{ij}$ se definen a
través de la derivada de la base coordenada,
\[
\frac{\partial \vec e_i}{\partial x^j} = \Gamma^k_{ij}\, \vec e_k .
\]
Miden el cambio de dirección (y de magnitud) de los vectores base al ser
trasladados a lo largo de curvas del espacio.
\end{definicion}

\begin{proposicionc}
Bajo la transformación de coordenadas $x'^i=x'^i(x^j)$, los símbolos de
Christoffel transforman como
\[
\Gamma'^{\,k}_{ij} = \frac{\partial x'^k}{\partial x^m}\frac{\partial
x^l}{\partial x'^i}\frac{\partial x^n}{\partial x'^j}\, \Gamma^m_{ln} +
\frac{\partial x'^k}{\partial x^m}\, \frac{\partial^2 x^m}{\partial
x'^i\partial x'^j} .
\]
\end{proposicionc}

\begin{observacion}
Los símbolos de Christoffel \textbf{no son un tensor}: su ley de
transformación contiene el término inhomogéneo
$(\partial x'^k/\partial x^m)(\partial^2 x^m/\partial x'^i\partial
x'^j)$, que depende del sistema de coordenadas, mientras que los
tensores, al ser covariantes, no dependen de él. Si los símbolos de
Christoffel fueran un tensor, transformarían únicamente con el primer
término (homogéneo),
$\Gamma'^{\,k}_{ij}=(\partial x'^k/\partial x^m)(\partial x^l/\partial
x'^i)(\partial x^n/\partial x'^j)\Gamma^m_{ln}$.

El primer término (homogéneo) mide el cambio \emph{covariante} de los
vectores base al ser trasladados, lo cual no depende del sistema
coordenado sino del tipo de espacio (curvo o plano) en el que se
trasladan. El segundo término (inhomogéneo) mide el cambio de los
vectores base al ser trasladados sobre líneas coordenadas curvas, y sí
depende del sistema coordenado elegido.
\end{observacion}

\section{Operaciones con tensores}

Los tensores viven en espacios vectoriales, por lo que satisfacen las
operaciones propias de un espacio vectorial.

\begin{definicion}[Suma]
Sean $V^{i_1\dots i_p}_{\ \ j_1\dots j_q}$ y $W^{i_1\dots
i_p}_{\ \ j_1\dots j_q}$ dos tensores del mismo orden $(p,q)$ (por
homogeneidad, deben serlo). Se define la suma
\[
T^{i_1\dots i_p}_{\quad\ j_1\dots j_q} = \alpha\, V^{i_1\dots
i_p}_{\quad\ j_1\dots j_q} + \beta\, W^{i_1\dots i_p}_{\quad\ j_1\dots
j_q}, \qquad \alpha,\beta\in K.
\]
\end{definicion}

\begin{definicion}[Producto tensorial]
El producto tensorial entre un tensor $V^{i_1\dots
i_p}_{\quad\ j_1\dots j_q}$ de orden $(p,q)$ y un tensor $W^{i_1\dots
i_{p'}}_{\quad\ j_1\dots j_{q'}}$ de orden $(p',q')$ se define como
\[
M^{i_1\dots i_{p+p'}}_{\quad\quad\ \ j_1\dots j_{q+q'}} = V^{i_1\dots
i_p}_{\quad\ j_1\dots j_q}\, W^{i_{p+1}\dots i_{p+p'}}_{\quad\quad\
j_{q+1}\dots j_{q+q'}} .
\]
\end{definicion}

\begin{proposicionc}
El producto tensorial de un tensor de orden $(p,q)$ con uno de orden
$(p',q')$ es un tensor de orden $(p+p',q+q')$.
\end{proposicionc}

Por lo tanto, el producto tensorial permite \textbf{construir tensores de
mayor orden}. La operación complementaria, llamada \textbf{contracción},
permite construir tensores de \textbf{menor orden}.

\begin{ejemplo}
Sea el tensor $K^i_{\ jl}$. Sumando (contrayendo) los índices $i$ y $l$,
\[
K^l_{\ jl} \qquad \text{(convención de Einstein: suma sobre } l\text{)},
\]
se obtiene una cantidad con un solo índice libre, $j$.
\end{ejemplo}

\begin{proposicionc}
Sea el tensor $K^i_{\ jl}$ de orden $(1,2)$; la contracción $K^l_{\ jl}$
es un tensor de orden $(0,1)$.
\end{proposicionc}

\begin{definicion}[Contracción]
La \textbf{contracción} consiste en poner bajo la operación de suma la
misma cantidad de índices contravariantes y covariantes.
\end{definicion}

\begin{proposicionc}
Sea $V^{i_1\dots i_p}_{\quad\ j_1\dots j_q}$ un tensor de orden $(p,q)$.
Efectuando una contracción de $r$ índices contravariantes con $r$
índices covariantes ($r<\min\{p,q\}$),
\[
V^{i_1\dots i_s k_1\dots k_r\, i_{s+r+1}\dots i_p}_{\quad\quad\quad\quad\
\ \, j_1\dots j_s k_1\dots k_r\, j_{s+r+1}\dots j_q},
\]
se obtiene un tensor de orden $(p-r,q-r)$.
\end{proposicionc}

\begin{observacion}
Si se desea contraer dos índices de la misma naturaleza (por ejemplo, dos
índices covariantes $j,k$ del tensor $T_{ijk}$), esto no cumple la
definición de contracción directamente. La solución es subir uno de los
dos índices con la métrica inversa antes de contraer:
\[
T_i^{\ k} = g^{kl}T_{ijl}\, \delta^j_{(\cdot)} \quad\longrightarrow\quad
\text{primero } T^{\ \ k}_{ij} = g^{kl}T_{ijl}, \quad\text{luego se
contrae } i \text{ con } k.
\]
\end{observacion}

\section{Simetrías}

\begin{definicion}[Tensor simétrico]
Un tensor se dice \textbf{simétrico} si bajo el intercambio de dos
índices el tensor no cambia:
\[
T^{\dots i\dots j\dots} = +T^{\dots j\dots i\dots}, \qquad T_{\dots
i\dots j\dots} = +T_{\dots j\dots i\dots}.
\]
\end{definicion}

\begin{ejemplo}
La métrica y la métrica inversa son simétricas: $g_{ij}=g_{ji}$,
$g^{ij}=g^{ji}$.
\end{ejemplo}

\begin{observacion}
La simetría de un tensor afecta el número de componentes independientes.
Por ejemplo, la métrica $g_{ij}$ posee, en principio, $n\times n=n^2$
componentes, pero al ser simétrica solo tiene $n(n+1)/2$ componentes
independientes.
\end{observacion}

\begin{proposicionc}
Un tensor de orden $(p,0)$, $T^{i_1\dots i_p}$, posee en principio $n^p$
componentes independientes. Si dos de sus índices son simétricos,
\[
T^{i_1\dots j\dots k\dots i_p} = +T^{i_1\dots k\dots j\dots i_p},
\]
el número de componentes independientes se reduce a
\[
n^{p-2}\times \frac{n(n+1)}{2} = \frac{1}{2}n^{p-1}(n+1).
\]
\end{proposicionc}

La misma conclusión se obtiene para tensores de orden $(0,q)$.

\begin{definicion}[Tensor antisimétrico]
Un tensor se dice \textbf{antisimétrico} si bajo el intercambio de dos
índices el tensor cambia de signo:
\[
T^{\dots i\dots j\dots} = -T^{\dots j\dots i\dots}, \qquad T_{\dots
i\dots j\dots} = -T_{\dots j\dots i\dots}.
\]
\end{definicion}

\begin{proposicionc}
Un tensor de orden $(2,0)$ antisimétrico, $T^{ij}=-T^{ji}$, en un espacio
plano de $n$ dimensiones posee $n(n-1)/2$ componentes independientes (en
vez de $n^2$). En general, para un tensor de orden $(p,0)$ con dos
índices antisimétricos,
\[
n^{p-2}\times\frac{n(n-1)}{2} = \frac{1}{2}n^{p-1}(n-1).
\]
La misma conclusión se obtiene para tensores de orden $(0,q)$.
\end{proposicionc}

\begin{proposicionc}
Un tensor antisimétrico posee ceros en los índices que se repiten:
$T^{\dots i\dots i\dots}=0$ (sin suma).
\end{proposicionc}

\begin{proof}
Haciendo $i=j$ en la definición de antisimetría (sin sumar sobre $i$),
$T^{\dots i\dots i\dots} = -T^{\dots i\dots i\dots}$, de donde
$T^{\dots i\dots i\dots}=0$. \qedhere
\end{proof}

\begin{definicion}[Descomposición simétrica-antisimétrica]
Dado un tensor $T^{ij}$, siempre es posible construir su parte simétrica
y antisimétrica:
\[
T^{(ij)} := \frac{1}{2}\left(T^{ij}+T^{ji}\right) = T^{(ji)}, \qquad
T^{[ij]} := \frac{1}{2}\left(T^{ij}-T^{ji}\right) = -T^{[ji]} .
\]
El factor $1/2$ evita contar dos veces. Lo mismo se aplica a $T_{ij}$.
\end{definicion}

\begin{proposicionc}
Todo tensor $T^{ij}$ (o $T_{ij}$) se descompone como
\[
T^{ij} = T^{(ij)} + T^{[ij]} .
\]
\end{proposicionc}

\begin{proposicionc}
Si $A^{ij}$ es simétrico y $B_{ij}$ es antisimétrico, entonces
$A^{ij}B_{ij}=0$.
\end{proposicionc}

\begin{proposicionc}
Si $A^{ij}$ es antisimétrico y $B_{ij}$ es un tensor cualquiera, entonces
$A^{ij}B_{ij} = A^{ij}B_{[ij]}$.
\end{proposicionc}

\begin{ejemplo}
Sea $T^{jkl}_{\ \ \ i}$; formemos sus partes simétrica y antisimétrica en
los índices $i,k$. Como $i$ es covariante y $k$ contravariante, son de
distinta naturaleza y deben igualarse antes de (anti)simetrizar. Subiendo
$i$ con la métrica inversa, $T^{ijkl}=g^{in}T^{jkl}_{\ \ \ n}$:
\begin{align*}
T^{(i|j|k)l} &= \frac{1}{2}\left(T^{ijkl}+T^{kjil}\right) =
\frac{1}{2}\left(g^{in}T^{jkl}_{\ \ \ n} + g^{kn}T^{jil}_{\ \ \ n}\right)
= g^{n(i|}T^{j|k)l}_{\quad\ n}, \\
T^{[i|j|k]l} &= \frac{1}{2}\left(T^{ijkl}-T^{kjil}\right) = g^{n[i|}
T^{j|k]l}_{\quad\ n},
\end{align*}
donde las barras verticales indican que el índice $j$ queda excluido de
la (anti)simetrización. Alternativamente, bajando el índice $k$ con la
métrica $g_{ij}$ se obtiene la descomposición equivalente
$T^{j\ \ l}_{\ (i\,k)}$, $T^{j\ \ l}_{\ [i\,k]}$.
\end{ejemplo}

\begin{definicion}[Tensor completamente antisimétrico]
Un tensor se dice \textbf{completamente antisimétrico} si cambia de
signo bajo cualquier permutación de un par de índices.
\end{definicion}

El caso completamente antisimétrico es de particular interés (más que el
completamente simétrico). Para un tensor de orden $(0,q)$ completamente
antisimétrico, $T_{i_1i_2\dots i_q}$, el número de componentes
independientes es
\[
\binom{n}{q} = \frac{n!}{q!(n-q)!} .
\]

\begin{observacion}
Si $q>n$ (más índices que dimensiones), entonces $\binom{n}{q}=0$: el
tensor es idénticamente nulo, pues los índices deben repetirse
obligatoriamente y, como vimos, los índices repetidos anulan un tensor
completamente antisimétrico.
\end{observacion}

Si $q=n$, el tensor completamente antisimétrico posee $\binom{n}{n}=1$
componente independiente.

\subsection{El símbolo de Levi-Civita}

\begin{ejemplo}[Caso $n=2$]
El tensor $T_{ij}$ completamente antisimétrico en $2$ dimensiones,
$T_{ij}=-T_{ji}$, se escribe
\[
T_{ij} = \begin{pmatrix} 0 & T_{12} \\ -T_{12} & 0 \end{pmatrix} =
T_{12} \begin{pmatrix} 0 & 1 \\ -1 & 0 \end{pmatrix} = T_{12}\,
\epsilon_{ij},
\]
posee una sola componente independiente $T_{12}$. Se define
\[
\epsilon_{ij} = \begin{cases} +1, & (ij)=(12) \\ -1, & (ij)=(21) \\ 0, &
i=j \end{cases}
\]
llamado \textbf{símbolo de Levi-Civita} en dos dimensiones.
\end{ejemplo}

\begin{observacion}
El símbolo de Levi-Civita \textbf{no} es un tensor de orden $(0,2)$, sino
una \emph{densidad tensorial} de orden $(0,2)$ y peso $-1$.
\end{observacion}

\begin{ejemplo}[Caso $n=3$]
El tensor totalmente antisimétrico $T_{ijk}$ en tres dimensiones posee
una única componente independiente $T_{123}$, tal que
\[
T_{ijk} = T_{123}\, \epsilon_{ijk},
\]
donde $\epsilon_{ijk}$ es el símbolo de Levi-Civita en tres dimensiones,
\[
\epsilon_{ijk} = \begin{cases} +1, & (ijk) \text{ permutación par de }
(123) \\ -1, & (ijk) \text{ permutación impar de } (123) \\ 0, & \text{en
otro caso.} \end{cases}
\]
Por ejemplo, $\epsilon_{213}=-\epsilon_{123}=-1$,
$\epsilon_{132}=-\epsilon_{123}=-1$, $\epsilon_{312} =
-\epsilon_{132}=+1$. Los signos, usando índices simbólicos $ijk$ en vez
de $123$, se obtienen por \textbf{permutación cíclica}: por ejemplo,
$\epsilon_{jik}=-\epsilon_{ijk}$, $\epsilon_{ikj}=-\epsilon_{ijk}$,
$\epsilon_{kji} = -\epsilon_{jki} = \epsilon_{jik}=-\epsilon_{ijk}$.
\end{ejemplo}

Las propiedades siguientes se establecen en un sistema coordenado
cartesiano de un espacio plano, donde $g_{ij}=\delta_{ij}$.

\begin{proposicionc}
El símbolo de Levi-Civita $\epsilon_{ijk}$ se puede escribir como
\[
\epsilon_{ijk} = \det\begin{pmatrix} \delta_{i1} & \delta_{i2} &
\delta_{i3} \\ \delta_{j1} & \delta_{j2} & \delta_{j3} \\ \delta_{k1} &
\delta_{k2} & \delta_{k3} \end{pmatrix},
\]
y satisface
\[
\epsilon_{ijk}\epsilon_{lmn} = \det\begin{pmatrix} \delta_{il} &
\delta_{im} & \delta_{in} \\ \delta_{jl} & \delta_{jm} & \delta_{jn} \\
\delta_{kl} & \delta_{km} & \delta_{kn} \end{pmatrix}, \qquad
\epsilon_{ijk}\epsilon_{lmk} = \delta_{il}\delta_{jm} -
\delta_{im}\delta_{jl}, \qquad \epsilon_{ijk}\epsilon_{ljk} =
2\delta_{il} .
\]
\end{proposicionc}

\begin{proposicionc}
El determinante de un tensor $A^{ij}$ se puede escribir como
\[
\det(A^{ij}) = \epsilon_{ijk}A^{i1}A^{j2}A^{k3} = \frac{1}{3!}\,
\epsilon_{ijk}\epsilon_{lmn}\, A^{il}A^{jm}A^{kn} .
\]
\end{proposicionc}

\section{Producto vectorial}

La base canónica $\{\ihat,\jhat,\khat\}$ de $\mathbb{R}^3$ satisface
$\ihat\times\jhat=\khat$, $\ihat\times\khat=-\jhat$,
$\jhat\times\khat=\ihat$. Escribiendo $\vec e_1=\ihat$, $\vec
e_2=\jhat$, $\vec e_3=\khat$, esta álgebra se resume como
\[
\vec e_i \times \vec e_j = \epsilon_{ijk}\, \vec e_k .
\]

\begin{proposicionc}
El producto cruz entre $\vec A$ y $\vec B$, $\vec C = \vec A\times\vec
B$, se escribe en notación de índices como
\[
C^i = \epsilon^{ijk}A_jB_k .
\]
\end{proposicionc}

\begin{proof}
\[
\vec C = \vec A\times\vec B = A^i\vec e_i\times B^j\vec e_j = A^iB^j\,
\vec e_i\times\vec e_j = A^iB^j\epsilon_{ijk}\vec e_k = \epsilon_{jki}
A^jB^k \vec e_i = \epsilon_{ijk}A^jB^k\, \vec e_i,
\]
usando la simetría cíclica $\epsilon_{jki}=\epsilon_{ijk}$. Comparando
con $\vec C = C^i\vec e_i$, se concluye $C^i = \epsilon^{ijk}A_jB_k$
(en cartesianas, $\epsilon^{ijk}=\epsilon_{ijk}$ e índices arriba/abajo
son intercambiables). \qedhere
\end{proof}

\begin{ejemplo}
Desarrollando explícitamente para $i=1,2,3$ (usando
$\epsilon_{1jk},\epsilon_{2jk},\epsilon_{3jk}$ y anulando los términos
con índices repetidos) se recupera la fórmula usual del producto cruz:
\[
C^1 = A^2B^3-A^3B^2, \qquad C^2 = A^3B^1-A^1B^3, \qquad C^3 =
A^1B^2-A^2B^1.
\]
\end{ejemplo}

\begin{proposicionc}
El vector $\vec C=\vec A\times\vec B$ es perpendicular a $\vec A$ y a
$\vec B$.
\end{proposicionc}

\begin{proof}
Calculemos $\vec C\cdot\vec A = C^iA_i = \epsilon^{ijk}A_jB_kA_i =
\epsilon^{ijk}A_iA_jB_k$. Definiendo el tensor simétrico
$T^{ij}=A^iA^j=A^jA^i=T^{ji}$,
\[
\vec C\cdot\vec A = \epsilon^{ijk}T^{ij}B_k .
\]
Por ser $T^{ij}$ simétrico y $\epsilon^{ijk}$ antisimétrico en $(i,j)$,
esta contracción se anula (Proposición sobre $A^{ij}B_{ij}=0$ para $A$
simétrico y $B$ antisimétrico):
\[
\epsilon^{ijk}T^{ij}B_k = -\epsilon^{jik}T^{ij}B_k = -\epsilon^{ijk}
T^{ji}B_k = -\epsilon^{ijk}T^{ij}B_k \;\Rightarrow\; \vec C\cdot\vec A =
0 .
\]
La misma conclusión se obtiene para $\vec C\cdot\vec B=0$. \qedhere
\end{proof}

Se sigue que $\vec A\times\vec B$ es perpendicular al plano que
contienen $\vec A$ y $\vec B$, y por lo tanto todo vector contenido en
dicho plano es perpendicular a $\vec A\times\vec B$: en efecto, todo
vector $\vec D$ del plano generado por $\vec A$ y $\vec B$ se escribe
\[
\vec D = \frac{D}{\sin(\theta_1+\theta_2)}\left(\sin\theta_2\, \hat A +
\sin\theta_1\, \hat B\right),
\]
con $D=\|\vec D\|$, $\hat A=\vec A/A$, $\hat B=\vec B/B$, de donde
$\vec D\cdot\vec C = 0$ directamente.

\begin{proposicionc}
Si $\vec A,\vec B,\vec C$ están contenidos en un plano, entonces
$\epsilon_{ijk}A^iB^jC^k=0$. Si son no coplanares, el volumen que forman
es
\[
V = \epsilon_{ijk}A^iB^jC^k .
\]
\end{proposicionc}

\begin{proposicionc}
La norma de $\vec C=\vec A\times\vec B$ es
\[
\|\vec C\| = \|\vec A\|\,\|\vec B\|\,\sin\theta ,
\]
donde $\theta$ es el ángulo entre $\vec A$ y $\vec B$.
\end{proposicionc}

\begin{proof}
\begin{align*}
\|\vec C\|^2 &= C^iC_i = \epsilon^{ijk}A_jB_k\,\epsilon_{inm}A^nB^m \\
&= \left(\delta^j_n\delta^k_m - \delta^j_m\delta^k_n\right) A_jB_kA^nB^m
\\
&= A_jA^jB_kB^k - A_jB^jA^kB_k \\
&= \|\vec A\|^2\|\vec B\|^2 - (\vec A\cdot\vec B)^2,
\end{align*}
usando la identidad $\epsilon^{ijk}\epsilon_{inm} =
\delta^j_n\delta^k_m-\delta^j_m\delta^k_n$. Como $\vec A\cdot\vec B =
\|\vec A\|\|\vec B\|\cos\theta$,
\[
\|\vec C\|^2 = \|\vec A\|^2\|\vec B\|^2(1-\cos^2\theta) = \|\vec
A\|^2\|\vec B\|^2\sin^2\theta \;\Rightarrow\; \|\vec C\| = \|\vec
A\|\,\|\vec B\|\,\sin\theta . \qedhere
\]
\end{proof}

\section{Derivada covariante}

Sea $\vec A$ un vector escrito en la base coordenada $Ox^i$. Calculemos
su derivada parcial $\partial_i := \partial/\partial x^i$:
\[
\partial_i \vec A = \partial_i(A^j\vec e_j) = (\partial_iA^j)\vec e_j +
A^j(\partial_i\vec e_j) = (\partial_iA^j)\vec e_j + A^j\Gamma^k_{ij}\vec
e_k = \left(\partial_iA^j + \Gamma^j_{ik}A^k\right)\vec e_j ,
\]
donde se usó la definición de los símbolos de Christoffel y se
renombraron índices mudos.

\begin{definicion}[Derivada covariante de un vector]
Sea $A^i$ un tensor de orden $(1,0)$. Se define su \textbf{derivada
covariante} como
\[
\nabla_i A^j := \partial_i A^j + \Gamma^j_{ik}A^k .
\]
\end{definicion}

Con esta definición, $\partial_i\vec A = \nabla_iA^j\,\vec e_j$: la
derivada covariante $\nabla_iA^j$ reemplaza, en notación tensorial, a la
derivada parcial $\partial_iA^j$.

\begin{proposicionc}
Si $A^i$ es un tensor de orden $(1,0)$, entonces $\nabla_iA^j$ es un
tensor de orden $(1,1)$.
\end{proposicionc}

Análogamente, para un vector dual $\vec f\in V^*(n,K)$,
\[
\partial_i\vec f = \partial_i(f_jE^j) = (\partial_if_j)E^j +
f_j\partial_iE^j = (\partial_if_j)E^j - f_k\Gamma^k_{ij}E^j =
\left(\partial_if_j - \Gamma^k_{ij}f_k\right)E^j ,
\]
donde se usó $\partial_iE^j = -\Gamma^j_{ik}E^k$ (consecuencia de derivar
$E^i(e_j)=\delta^i_j$).

\begin{definicion}[Derivada covariante de un covector]
Sea $A_i$ un tensor de orden $(0,1)$. Se define
\[
\nabla_iA_j := \partial_iA_j - \Gamma^k_{ij}A_k .
\]
\end{definicion}

\begin{proposicionc}
Si $A_i$ es un tensor de orden $(0,1)$, entonces $\nabla_iA_j$ es un
tensor de orden $(0,2)$.
\end{proposicionc}

\begin{observacion}
Resumiendo: $\nabla_iA^j = \partial_iA^j+\Gamma^j_{ik}A^k$ (el índice
arriba aparece con $+\Gamma$); $\nabla_iA_j =
\partial_iA_j-\Gamma^k_{ij}A_k$ (el índice abajo aparece con $-\Gamma$).
\end{observacion}

Para un tensor mixto $A\in V(n,\mathbb{R})\otimes V^*(n,\mathbb{R})$, un
cálculo análogo (derivando $A=A^j_{\ k}\,\vec e_j\otimes E^k$ y usando
las leyes de derivación de ambas bases) da
\[
\partial_iA = \left(\partial_iA^j_{\ k} + \Gamma^j_{il}A^l_{\ k} -
\Gamma^l_{ik}A^j_{\ l}\right)\vec e_j\otimes E^k .
\]

\begin{definicion}[Derivada covariante de un tensor $(1,1)$]
Sea $A^i_{\ j}$ un tensor de orden $(1,1)$. Se define
\[
\nabla_iA^j_{\ k} = \partial_iA^j_{\ k} + \Gamma^j_{il}A^l_{\ k} -
\Gamma^l_{ik}A^j_{\ l} .
\]
\end{definicion}

\begin{proposicionc}
Si $A^i_{\ j}$ es un tensor de orden $(1,1)$, entonces $\nabla_iA^j_{\
k}$ es un tensor de orden $(1,2)$.
\end{proposicionc}

La generalización a tensores de orden $(p,q)$ arbitrario es directa:

\begin{definicion}[Derivada covariante general]
Sea $A^{i_1\dots i_p}_{\quad\ j_1\dots j_q}$ un tensor de orden $(p,q)$.
Se define su derivada covariante como
\[
\nabla_i A^{j_1\dots j_p}_{\quad\ k_1\dots k_q} = \partial_i A^{j_1\dots
j_p}_{\quad\ k_1\dots k_q} + \Gamma^{j_1}_{il}A^{lj_2\dots j_p}_{\quad\
\ k_1\dots k_q} + \dots + \Gamma^{j_p}_{il}A^{j_1\dots j_{p-1}l}_{\quad\
\quad\ \, k_1\dots k_q} - \Gamma^l_{ik_1}A^{j_1\dots j_p}_{\quad\
lk_2\dots k_q} - \dots - \Gamma^l_{ik_q}A^{j_1\dots j_p}_{\quad\
k_1\dots k_{q-1}l} .
\]
\end{definicion}

\begin{proposicionc}
Si $A^{i_1\dots i_p}_{\quad\ j_1\dots j_q}$ es un tensor de orden $(p,q)$,
entonces $\nabla_iA^{j_1\dots j_p}_{\quad\ k_1\dots k_q}$ es un tensor
de orden $(p,q+1)$.
\end{proposicionc}

\begin{proposicionc}[Compatibilidad de la métrica]
La métrica satisface
\[
\nabla_i g_{jk} = 0 .
\]
\end{proposicionc}

\section{El operador rotacional}

Trabajando en coordenadas cartesianas y usando notación de índices, el
operador rotacional se define como
\begin{align*}
\vec B = \vec\nabla\times\vec A &= E^i\partial_i\times(A^j\vec e_j) \\
&= \delta^{ik}\vec e_k \times \big[(\partial_iA^j)\vec e_j + A^j
\cancel{\partial_i\vec e_j}\big] \\
&= \vec e_k\times(\partial_iA^j)\vec e_j
= \partial_iA^j\,\vec e_i\times\vec e_j = \partial_iA^j\epsilon_{ijk}\vec
e_k ,
\end{align*}
donde el término $\partial_i\vec e_j$ se anula por trabajar en
coordenadas cartesianas ($\Gamma^k_{ij}=0$). Renombrando índices,
$\partial_jA^k\epsilon_{kij}\vec e_i = \epsilon_{ijk}\partial_jA^k\,\vec
e_i$, de donde, comparando con $\vec B = B^i\vec e_i$:

\begin{proposicionc}
En coordenadas cartesianas, el rotacional de $\vec A$ se escribe en
notación de índices como
\[
B^i = \epsilon^{ijk}\partial_jA^k .
\]
\end{proposicionc}

\section{Notación tensorial}

A partir de aquí adoptamos la \textbf{notación tensorial}: se usan los
tensores directamente, obviando los vectores base. La correspondencia es

\begin{center}
\begin{tabular}{ccc}
Notación vectorial & & Notación tensorial \\ \hline
$\vec x(x^i)$ & $\longrightarrow$ & $x^i$ \\
$\vec A(x^i)$ & $\longrightarrow$ & $A^i(x^j)$ \\
$\vec A\cdot\vec B$ & $\longrightarrow$ & $g_{ij}A^iB^j$ \\
$\vec\nabla$ & $\longrightarrow$ & $\partial/\partial x^i =: \partial_i$
\\
$\partial_iA^{j_1\dots j_p}_{\quad\ k_1\dots k_q}$ & $\longrightarrow$ &
$\nabla_iA^{j_1\dots j_p}_{\quad\ k_1\dots k_q}$
\end{tabular}
\end{center}

\begin{observacion}
En coordenadas cartesianas se recupera el caso particular
\[
\{\vec e_i\}_{i=1}^3 = \{\ihat,\jhat,\khat\}, \qquad g_{ij}=\delta_{ij},
\qquad \Gamma^k_{ij}=0, \qquad \nabla_i = \partial_i .
\]
\end{observacion}

\section{Invariancia de la acción}

Un sistema físico de $f$ grados de libertad tiene asociado un principio
de acción
\[
S = \int_{t_i}^{t_f} L\big(q^i,\dot q^i,t\big)\, dt ,
\]
del cual se deducen las ecuaciones de movimiento usando el operador
variacional,
\[
\frac{\delta L}{\delta q^i} = \frac{\partial L}{\partial q^i} -
\frac{d}{dt}\left(\frac{\partial L}{\partial \dot q^i}\right).
\]
El fundamento del principio de acción es que $S$ es un número real: se
mapea una trayectoria en el espacio de configuración a un número real.
Las ecuaciones de movimiento se obtienen exigiendo que $S$ tenga un
extremo sobre la curva que describe el movimiento físico del sistema.

Al cambiar de coordenadas, $q'^i=q'^i(q^j)$, el Lagrangiano cambiaría, en
principio, su forma funcional:
\[
S' = \int_{t_i}^{t_f} L'\big(q'^i,\dot q'^i,t\big)\, dt .
\]
Sin embargo, la curva que extremiza la acción no cambia su forma física;
solo cambian las ecuaciones (por el cambio de coordenadas) que la
describen. Por lo tanto, el valor numérico de la acción no cambia,
$S=S'$:
\[
\int_{t_i}^{t_f} L\big(q^i,\dot q^i,t\big)\, dt = \int_{t_i}^{t_f}
L'\big(q'^i,\dot q'^i,t\big)\, dt \quad\Longrightarrow\quad
L'\big(q'^i,\dot q'^i,t\big) = L\big(q^i,\dot q^i,t\big),
\]
que es la definición del nuevo Lagrangiano $L'$: el Lagrangiano es un
\textbf{escalar} bajo transformaciones de coordenadas.

\begin{ejemplo}
Consideremos el problema de fuerzas centrales conservativas, con
Lagrangiano en coordenadas cartesianas
\[
L(x,y,\dot x,\dot y) = \frac{1}{2}m(\dot x^2+\dot y^2) -
U\!\left(\sqrt{x^2+y^2}\right).
\]
Cambiando a coordenadas cilíndricas (con $z=0$), $x=r\cos\varphi$,
$y=r\sin\varphi$, se obtiene
\[
L'(r,\varphi,\dot r,\dot\varphi) = \frac{1}{2}m\left[\left(\frac{d}{dt}
(r\cos\varphi)\right)^2 + \left(\frac{d}{dt}(r\sin\varphi)\right)^2
\right] - U(r) = \frac{1}{2}m\left(\dot r^2 + r^2\dot\varphi^2\right) -
U(r).
\]
\end{ejemplo}

En una transformación de coordenadas no solo cambia la función
Lagrangiana, sino también la forma funcional de las ecuaciones de
movimiento.

\begin{proposicionc}
Bajo una transformación de coordenadas $q'^i=q'^i(q^j)$, las ecuaciones
de movimiento cambian como
\[
\frac{\delta L'}{\delta q'^i} = \frac{\partial q^j}{\partial q'^i}\,
\frac{\delta L}{\delta q^j} .
\]
Es decir, las ecuaciones de movimiento $\delta L/\delta q^i$ forman un
tensor de orden $(0,1)$.
\end{proposicionc}



\end{document}
